\section{Part D. Rational points behind sections}
\label{sec:partD}

The results below isolate actual implications and supply abstract
discontinuous actions. D.14 refutes the universal assertion D5;
D1--D4 remain unresolved. D.15--D.19 reconstruct and check the
local proof supplied as R24, and D.6 gives its affine consequence.
These local results do not supply the missing global hypotheses.
The manuscript passages are cote 149, manuscript pp.~22--31,
scans 45, 48, 56, 59--63, and pp.~62v--63, scans 126--127.

\subsection{The lifting assertions and their quantifiers}

\begin{claim}{D.1}{Known; Bresciani's lifting criterion and its consequence}
For a smooth geometrically connected curve $X/k$, with $k$ finitely
generated over $\Q$, a continuous section is geometric or cuspidal
if and only if its base change to $k(t)$ is birationally liftable.
Consequently, the assertion D3, quantified over every such field and
every hyperbolic curve, is equivalent to the section conjecture with
cuspidal sections allowed at the rational boundary.
\end{claim}
\begin{proof}
The first assertion is precisely Theorem A of the opened final source
\citet{D_Bresciani2023}. Its additional hypothesis is liftability
\emph{after the transcendental extension} $k(t)/k$; liftability only
over $k$ is not substituted for it.
If D3 holds, apply it to $X_{k(t)}$ and the restriction of any given
section of $X/k$. Theorem A makes that original section geometric or
cuspidal. Conversely a geometric section at $x\in X(k)$ lifts through
the decomposition group of the discrete valuation of $k(X)$ at $x$.
In characteristic zero its exact sequence has kernel
$\widehat{\Z}(1)$ and is split after choosing a local parameter:
the coefficient action on the Puiseux extensions gives a continuous
splitting. A cuspidal section already factors through that same
decomposition group at the boundary, so also lifts. Composing with
$G_{k(X)}\twoheadrightarrow\pi_1(X)$ gives the required original
section up to geometric conjugacy; conjugating the lift adjusts the
representative. These maps are continuous maps of profinite groups
over $G_k$. This is also the derivation stated in Bresciani's
Theorem B, with exactly its universal quantifier over finitely
generated extensions.
\end{proof}

\begin{claim}{D.2}{Proved; D2 and D4 imply D1}
If D2 and D4 both hold, every continuous section of
$\pi_1(\mathbb P^1_\Q\setminus\{0,1,\infty\})\to G_\Q$ is a point
section or a cuspidal section. In particular one with $C(s)=1$ is a
point section, using the checked cyclotomic-centralizer assertion R15.
\end{claim}
\begin{proof}
Lift $s$ by D4 to $\widetilde s:G_\Q\to G_{\Q(t)}$.
By D2 its image belongs to $D_v$ for a rational place $v$ of
$\Q(t)$. If $v$ belongs to the open curve, the quotient
$G_{\Q(t)}\to\pi_1(U)$ kills its inertia. The image of $D_v$ is
the decomposition section at $v$, so the quotient of
$\widetilde s$ is $s_v$. If $v$ is one of the three removed
points, its inertia is retained and the image belongs to the
cuspidal decomposition group. Its nontrivial geometric inertia
centralizes the restriction to $G_\Q^\circ$, because the latter
has trivial cyclotomic character. Thus $C(s)=1$ excludes this case.
All choices of prolongation of $v$ change the conclusion by
conjugacy in $G_{\overline\Q(t)}$, hence by geometric conjugacy
after passage to $\pi_1(U)$.
\end{proof}

The converse direction required to settle D3 is not the established
Hom theorem for open homomorphisms: a section has no open image in
the relevant geometric extension. Nor does solving separate finite
embedding problems produce compatible solutions automatically.
For a countable base, successive surjectivity for all pairs of open
subcurves does supply a compatible sequence, as the lemma labelled
\texttt{cuspbl} in the opened Bresciani source proves. Surjectivity
of those transition maps is itself the missing assertion.

\subsection{The adelic integrality route}

\begin{claim}{D.3}{Proved; verification of the arithmetic part of R23}
Write $\widehat{\Q^\times}=\varprojlim_n\Q^\times/(\Q^\times)^n$.
Its elements are pairs
\[
 \left(\epsilon,(a_p)_p\right),\quad
 \epsilon\in\{\pm1\},\quad a_p\in\widehat{\Z},
 \quad\#\{p:a_p\not\equiv0\pmod n\}<\infty\quad(n\geq1).
\]
An element belongs to $\Q^\times$ if and only if every $a_p$ is an
ordinary integer and all but finitely many are zero. In particular,
if all its local Kummer classes are represented by elements of
$\Q_p^\times$, and almost all by units of $\Z_p$, it is rational.
Moreover
\[
 \varprojlim_n\mathbb A_\Q^\times/(\mathbb A_\Q^\times)^n
       \simeq\widehat{\Q^\times}\times\widehat{\Z}^{\times}.
\]
\end{claim}
\begin{proof}
Unique factorization gives
$\Q^\times=\{\pm1\}\times\bigoplus_p\Z$.
Reducing modulo $n$ gives finite-support exponent vectors; taking
the inverse limit gives exactly the displayed description. Conversely
such compatible vectors define an element at every level. The
valuation $v_p$ extends to the continuous homomorphism into
$\widehat\Z$ extracting $a_p$. A local representative has integral
valuation; a local unit has valuation zero. Finite integral support
and the sign give a rational number with the same coordinates, and
the coordinate description proves equality. This proves the stated
local criterion without a global approximation assumption.

An idele has finitely many nonzero finite valuations. Multiplying
by the unique rational number with those valuations and the sign of
its real component leaves a finite unit idele and a positive real
component. This gives the group isomorphism
$\mathbb A_\Q^\times\simeq
\Q^\times\times\widehat\Z^\times\times\mathbb R_{>0}$.
The last factor is divisible. For the compact abelian group
$B=\widehat\Z^\times$, the map $B\to\varprojlim B/B^n$ is
surjective by compactness. Its kernel is $\bigcap_n B^n=1$:
each finite quotient of each $\Z_\ell^\times$ is killed by its
order, and these quotients separate points. This proves the last
isomorphism. The local-to-global assertion agrees with the opened
Stix Proposition 2 \citep{D_Stix2015}; here the proof for $\Q$ is
written independently.
\end{proof}

\begin{claim}{D.4}{Refuted; naive evaluation at a zero or pole}
For a birational section, the pullback of the Kummer class of
$f\in\Q(t)^\times$ is not always the class of the value of $f$ at
its associated place. At a zero or pole it is a leading-coefficient
class multiplied by the chosen tangential class. This corrects the
unqualified formula in the proposed route to D2.
\end{claim}
\begin{proof}
At $t=0$ choose the splitting of
$G_{\Q((t))}\to G_\Q$ obtained by fixing all compatible roots of
$t$, and twist it by the Kummer class of $2$. Embed its decomposition
group into $G_{\Q(t)}$ by choosing a prolongation of the $t$-adic
valuation. For $f=t^m u(t)$, with $u(0)\ne0$, its Kummer pullback is
\[
                         \kappa\bigl(2^m u(0)\bigr).
\]
To check this, choose an $n$th root of $u(0)$ and apply Hensel's
lemma to $u(t)/u(0)$, whose constant term is one; its $n$th root
with constant term one is fixed by the coefficient action. The
root of $t^m$ contributes the twist $2^m$. This works compatibly
for every $n$ in characteristic zero. In particular the three
values for $t,t-1,t+1$ are $2,-1,1$. They cannot be values of a
single field evaluation: evaluation with $t=2$ would give $t-1=1$.
The place is nevertheless the rational place $t=0$. The identities
are checked symbolically in \texttt{code/root\_checks.py}.
\end{proof}

\begin{claim}{D.5}{Known; the exact sufficient integrality condition}
Let $U=\mathbb P^1_\Q\setminus\{0,1,\infty\}$. A continuous section
which is locally a point of $U(\Q_p)$ at every finite place and is
integral on $\mathbb P^1_\Z\setminus\{0,1,\infty\}$ outside one
finite set of primes is the section of a point of $U(\Q)$.
This does not supply the missing uniform finite set for D1 or D2.
\end{claim}
\begin{proof}
For precision, Stix's Corollary 6 in the opened final arXiv source
\citep{D_Stix2015} applies to a hyperbolic curve over $\Q$ or an
imaginary quadratic number field with a nonconstant global unit,
and identifies adelic sections with rational points. Here the unit
$t$ supplies the hypothesis. We describe the proof's arithmetic
step: the images of $s$ under $t$ and $1-t$ are classes
$a,b\in\widehat{\Q^\times}$. The local points and the uniform
integrality hypothesis put them in $\Q^\times$ by D.3. Injectivity
of local Kummer theory gives $x_p=a$ and $1-x_p=b$ for every $p$,
so $a+b=1$ and the local support is the singleton $\{a\}$.
The source's Corollary 5 converts finite support of a noncuspidal
Selmer section into a point section. Its proof applies Stoll's
Theorem 8.2 to the finite fibres over $\{a\}$ in all neighbourhood
covers, after a cover with compactification of positive genus,
and then takes their inverse limit. The hypothesis on finite
descent is supplied for section-induced local points by
Harari--Stix, Theorem~2.1 and Remark~2.2(1) in the attached
preprint (Theorem~11 in the published numbering quoted by Stix)
\citep{D_HarariStix2012,D_Stoll2007}. Its exceptional set of
places may have Dirichlet density zero, and its condition concerns
all torsors under finite group schemes. Both hypotheses have the
required form for a globally induced section. We also opened the
2017 errata to Stoll: the corrections concerning complex places
and connected torsors do not change this use of Theorem~8.2
\citep{D_StollErrata2017}.
The rational point at the limit projects to $a$, and its
decomposition section is $s$.

Stix's Theorem 2 gives, for a birationally liftable noncuspidal
section, integrality at a set of primes of Dirichlet density one.
That set need not be cofinite. For example any sequence of primes
growing faster than $\exp(2^n)$ is infinite and has Dirichlet
density zero. Thus that theorem does not supply the hypothesis
used in the first paragraph. The missing passage from density
one to a common finite exceptional set is retained as an
unresolved step, not made by changing the meaning of ``almost all''.
\end{proof}

\subsection{The affine consequence of the checked local proof}

\begin{claim}{D.6}{Proved}
For every finite extension $k/\Q_p$ and every hyperbolic curve
$U/k$, every continuous section of the full arithmetic
\'{e}tale fundamental group is the section of a rational point
of $U$, or is contained in the decomposition group of a rational
cusp. The proper local theorem used here is proved in D.15--D.19
by reconstructing the supplied R24 argument.
\end{claim}
\begin{proof}
Fix $k/\Q_p$ finite and a section $s$ of $U/k$. In its system of
finite \'etale neighbourhoods, the covers whose compactification
has genus at least two form a cofinal system. Indeed a geometric
hyperbolic curve has such a cover: for genus zero use a sufficiently
large Fermat cover after a three-puncture map; for a once-punctured
genus-one curve use a finite quotient of its free fundamental group
in which the boundary commutator is nontrivial, then
Riemann--Hurwitz; genus at least two needs no initial cover.
A characteristic refinement of the normal core descends to an
$s$-neighbourhood over $k$. Further covers preserve the
genus bound by Riemann--Hurwitz.

More explicitly, for an open characteristic subgroup
$M\subset\Pi_U$, the arithmetic subgroup $M\,s(G_k)$ defines
the neighbourhood over $k$. Its geometric curve is the quotient
of the universal geometric cover by $M$, with the semilinear
action induced by $s$. Characteristic refinements are cofinal
because $\Pi_U$ is topologically finitely generated. Their
compactifications inherit this action and the transition maps.

For each such neighbourhood $V$, the section on $V$ induces a
section on $\overline V$, hence a point of $\overline V(k)$ by
D.19. These are nonempty compact spaces in the
$p$-adic topology. Their continuous transition maps have a
nonempty inverse limit: every finite set of compatibility
conditions is realized in a common dominating neighbourhood, and
compactness applies to the product. The resulting completed
pro-point is fixed by $s(G_k)$. Its projection is either an
interior point, whose stabilizer is its decomposition section,
or a boundary point, whose stabilizer is its cusp decomposition
group. This is exactly the fixed-point dictionary of the
fundamental covering tower. A fixed pro-point projects to a
$k$-rational point of the smooth compactification, so the cusp
in this conclusion is rational over the original field, not
only over an unspecified finite extension.
\end{proof}

\subsection{Abstract semilinear actions need not be continuous}

\begin{claim}{D.7}{Proved; an explicit discontinuous family of actions}
There exist abstract semilinear actions of $G_\Q$ on the universal
pro-cover of $U=\mathbb P^1_{\overline\Q}\setminus\{0,1,\infty\}$
for which a rational function has a nonopen stabilizer. The examples
constructed here normalize the original geometric deck group and
have exactly one $G_\Q^\circ$-fixed completed pro-point. They
therefore refute automatic continuity, but do not refute D5.
\end{claim}
\begin{proof}
Use the rational base point $x=1/2$. Let $\rho_x$ be its continuous
semilinear action on the pointed universal pro-cover, and let
$h(t)=1-t$. Because $h(x)=x$, pullback of pointed covers gives a
unique lift $\widetilde h$ fixing the distinguished pro-point
$\widetilde x$. It has order two and commutes with $\rho_x(G_\Q)$:
both facts follow from the corresponding identities on the
pointed covering functor. It normalizes the deck group.

For distinct rational primes $p_1,p_2,\ldots$, the multiquadratic
extension has Galois group $V=\prod_{n\geq1}\F_2$. Independence
of the prime square classes proves that the continuous restriction
map $G_\Q\twoheadrightarrow V$ is surjective. Regard $V$ as an
abstract vector space. The constant-one vector does not belong to
$V_0=\bigoplus_n\F_2$. Choose a linear functional
$\lambda:V/V_0\to\F_2$ taking its class to one, and compose
with the quotient map. Such a functional exists by extending a
basis. It is nonzero but vanishes on the dense subgroup $V_0$,
so is discontinuous. Let $\epsilon:G_\Q\to\F_2$ be the composite.
It is discontinuous: a continuous surjection of compact groups is
a quotient map. It vanishes on $G_\Q^\circ$, since every quadratic
extension of $\Q$ is cyclotomic.

Define
\[
       \rho(g)=\widetilde h^{\epsilon(g)}\rho_x(g).
\]
Commutation and additivity of $\epsilon$ prove the action law.
On constants the action is still $g$, so it is semilinear.
The stabilizer of the function $t$ is exactly $\ker\epsilon$,
which is nonopen. Every $\rho(g)$ normalizes the original deck
group. On $G_\Q^\circ$ the action equals $\rho_x$, and fixes
$\widetilde x$.

We verify uniqueness rather than assume it from an earlier AI
return. Ribet's Theorem 1 in the opened appendix to Katz--Lang
\citep{D_Ribet1981} says that the torsion of an abelian variety over
a number-field cyclotomic extension is finite. It also implies
that no nonzero point is divisible by every integer: choose
compatible division points of such a point $P$ and descend $P$
and the variety to a number field $k_0$. Their Kummer cocycle on
$G_{k_0}$ vanishes on the normal cyclotomic subgroup. The cocycle
identity makes its values invariant under that subgroup, whereas
the invariant Tate module is zero by Ribet. Thus all division
points are already over $k_0$. Embed them in $A(k_{0,v})$ at a
finite place of good reduction. This is a profinite group; its
finite quotients include reduction modulo every power of the
maximal ideal. A point divisible by every integer maps to zero
in each finite quotient, and these quotients separate points.
Thus $P=0$; reduction to the residue field alone would not suffice.
This holds for every finite extension of $\Q^{\mathrm{ab}}$,
which is $k_0(\mu_\infty)$ for some number field $k_0$.

Here a completed pro-point means a compatible system of geometric
closed points on the proper finite levels; it does not mean the
generic scheme point. If a second such pro-point were fixed, choose a finite
characteristic covering level separating the two points and
dominating a cover with compactification of positive genus.
The subgroup generated by that characteristic geometric group
and $\rho_x(G_{\Q^{\mathrm{ab}}})$ defines a cover over
$\Q^{\mathrm{ab}}$. The two distinct projected points on its
compactification have the same section class. Their Abel--Jacobi
difference has zero Kummer class and is divisible by every integer,
hence zero by the preceding argument. A principal divisor of the
form $P-Q$ with $P\ne Q$ would give a degree-one map to
$\mathbb P^1$, contradicting positive genus. This proves the
uniqueness assertion.
\end{proof}

\begin{claim}{D.8}{Proved; a failed proposed route}
Reparameterizing the coefficient action by a nonidentity abstract
automorphism of $G_\Q$ does not give an action semilinear for the
specified action of $G_\Q$ on $\overline\Q$.
\end{claim}
\begin{proof}
If the coefficient action of $g$ is $\alpha(g)$, semilinearity
requires $\alpha(g)(a)=g(a)$ for every $a\in\overline\Q$.
The action defining $G_\Q$ is faithful, so $\alpha(g)=g$ for
every $g$. An individual field automorphism of $\overline\Q$
only conjugates the Galois action and induces a continuous inner
automorphism of its Krull topological group. It does not produce
the discontinuous cocycle demanded in the proposed route.
\end{proof}

\begin{claim}{D.9}{Proved; abelianized splitting is insufficient}
There is a finite extension of groups with perfect kernel which
has no section. Its quotient obtained by abelianizing the kernel
has a section. Therefore vanishing of the abelianized obstruction
alone is not a valid general solution of the embedding problems
in the proposed route to D3--D4.
\end{claim}
\begin{proof}
Put $F=\F_9$ and let $\sigma:x\mapsto x^3$ be its nontrivial
automorphism. In $\operatorname{PGL}_2(F)\rtimes\langle\sigma\rangle$
take the subgroup $G$ consisting of $\operatorname{PSL}_2(F)$
and the elements $[A]\sigma$ with nonsquare determinant.
The determinant square class is multiplicative and invariant
under $\sigma$, so this union is a subgroup. Its quotient by
$N=\operatorname{PSL}_2(F)$ is $\Z/2$.

If an element $[A]\sigma$ in the nontrivial coset had square one,
then $A A^\sigma=\lambda I$. Applying $\sigma$ and conjugating
by $A$ gives $\lambda^\sigma=\lambda$, so
$\lambda\in\F_3^*$. Taking determinants gives
\[
              (\det A)^4=\lambda^2=1.
\]
A nonsquare in the cyclic group $\F_9^*$ has fourth power
$-1$, a contradiction. Thus the nontrivial coset contains no
involution and the quotient has no section.

For completeness $N$ is perfect. The upper and lower unipotent
matrices generate $\operatorname{SL}_2(F)$ by elementary row
operations. Conjugation by $\operatorname{diag}(d,d^{-1})$,
with $d^2\ne1$, multiplies the upper parameter by $d^2$.
Its commutators with upper unipotents therefore give every upper
unipotent, since $d^2-1\ne0$. The analogous lower calculation
does the same for lower unipotents. Passing to the projective
quotient proves perfectness. Hence $G/[N,N]=G/N$ and the
abelianized-kernel extension is the identity on $\Z/2$.

The program \texttt{code/D\_embedding\_obstruction.py}
enumerates all $720$ projective matrices, checks the $360$
nonsquare-coset elements, and verifies that none squares to the
identity after the field involution. It also verifies generation
of the $360$-element kernel by elementary matrices. This is a
counterexample to the abstract inference, not to D3 or D4 for
absolute Galois groups: the latter require their particular
nonabelian kernels and local conditions to be used.
\end{proof}

% Extension prepared by the Part F agent and independently reviewed by root.
% Do not infer the universal fixed-point assertion D5 from the example below.
\subsection{An abstract action outside every finite quotient}

\begin{claim}{D.10}{Proved}
Let $\mathcal U$ be the universal pro-cover of
$U=\mathbb P^1_{\overline\Q}\setminus\{0,1,\infty\}$ and
$\Omega$ its function field. There are a continuous tangential
semilinear action $\rho_0$ of $G_\Q$, a completed cusp $c$, and
an injective homomorphism
\[
 T:(\mathbb A_f,+)\longrightarrow\operatorname{Aut}_{\overline\Q}(\mathcal U)
\]
such that every $T_a$ fixes $c$,
\[
 \rho_0(g)T_a\rho_0(g)^{-1}=T_{\chi(g)a},
 \qquad
 \operatorname{Stab}_{\mathbb A_f}(\lambda)=\widehat\Z.
\]
Here $\lambda$ is the coordinate on $U$, viewed in $\Omega$,
and $\mathbb A_f$ denotes the additive finite adeles.
\end{claim}
\begin{proof}
Fix $\overline\Q\hookrightarrow\C$. The classical modular
uniformization identifies
$\mathfrak H/\Gamma(2)$ with $U(\C)$, and for
$q=\exp(\pi i\tau)$ its coordinate has expansion
\[
 \lambda(q)=16q-128q^2+704q^3+\cdots\in\Q[[q]].
\]
The product formula and this uniformization are given in
\citet[p.~628, equation~(1.0.2) and the following paragraph]{F_CDT2025}.
Only this classical coordinate description is used, not that
paper's unbounded-denominators theorem. Analytic finite-cover
comparison is SGA~1, XII, Theorem~5.1 \citep{F_SGA1}.
Choose the compatible cusp above $\tau=i\infty$ to embed
\[
 \Omega\hookrightarrow
 \mathcal P:=\bigcup_{n\ge1}\overline\Q((q^{1/n})).
\]
An element of $\Omega$ lifts to a meromorphic function on
$\mathfrak H$, invariant under a finite-index subgroup of
$\Gamma(2)$ and meromorphic at its cusps. Its selected Puiseux
germ has algebraic coefficients. The coordinate ring of
$\mathcal U$ consists of the elements with no poles on
$\mathfrak H$; this will also check the affine pro-curve, not
just its field.

Every finite adele is $a=r+z$ with $r\in\Q$ and
$z\in\widehat\Z$, and $\Q\cap\widehat\Z=\Z$.
Thus the quotient $\mathbb A_f/\widehat\Z\simeq\Q/\Z$
defines the character
\[
 \psi:\mathbb A_f\longrightarrow\mu_\infty(\C),
 \qquad \psi(r+z)=\exp(2\pi i r).
\]
For a Puiseux series set
\[
 T_a\left(\sum_k b_kq^{k/n}\right)
       =\sum_k b_k\psi(ka/n)q^{k/n}.
\]
Changing a denominator $n$ to $mn$ changes $k$ to $mk$,
so this formula is well-defined. It gives
$T_aT_b=T_{a+b}$, fixes constants, and preserves the
$q$-valuation.

If $z\in\widehat\Z$, then $T_z$ fixes $q$ and hence
$\lambda(q)$. It is the usual cusp-inertia action on the
Puiseux roots. Since $\Omega/\overline\Q(\lambda)$ is normal
with allowed branch locus $\{0,1,\infty\}$, this preserves
$\Omega$ and the affine coordinate ring.
For $r\in\Q$, the operator $T_r$ is analytically
$f(\tau)\mapsto f(\tau+2r)$. The matrix
$u(2r)=\left(\begin{smallmatrix}1&2r\\0&1\end{smallmatrix}\right)$
commensurates $\Gamma(2)$: clearing a denominator $d$ shows that
conjugation of $\Gamma(2d^2)$ lands in $\Gamma(2)$, and the
same argument applies to the inverse. Therefore if $f$ is
invariant under $\Gamma_f$ of finite index in $\Gamma(2)$,
then $T_rf$ is invariant under the finite-index subgroup
\[
 \Gamma(2)\cap u(2r)^{-1}\Gamma_f u(2r).
\]
It is meromorphic at every cusp. Consequently $T_rf$ is
algebraic over $\C(\lambda)$ with branching only over
$0,1,\infty$.

We check descent here explicitly. The Puiseux coefficients of
$T_rf$ remain algebraic. Take a polynomial relation over
$\C(\lambda)$ and clear denominators. Its finitely many
monomials in $\lambda$ and $T_rf$ have Laurent coefficients in
$\overline\Q$. The homogeneous linear system obtained by
comparing coefficients has a nonzero complex solution, hence
a nonzero solution over $\overline\Q$: its rank is determined
by finite minors. Thus $T_rf$ is algebraic over
$\overline\Q(\lambda)$ and belongs to $\Omega$.
Moreover $\tau\mapsto\tau+2r$ preserves $\mathfrak H$ and
the rational cusps, so it preserves absence of poles on
$\mathfrak H$. The inverse is $T_{-r}$. Every $a=r+z$ now
acts on $\mathcal U$ as asserted. This proof does not assume
uniqueness of roots in the full automorphism group.

These automorphisms also act on the completed pro-cover.
Indeed the image of any finite-level function field is
contained in a further finite level. The resulting rational
map between the smooth proper compactifications extends
uniquely everywhere, by properness and normality for curves.
The extensions are compatible under composition and passage
to further levels. Applying the inverse operator proves that
they define an automorphism of the completed pro-object.

The coefficient of $q$ in $T_a\lambda$ is $16\psi(a)$.
It follows that $T_a\lambda=\lambda$ if and only if
$a\in\widehat\Z$: sufficiency follows because $T_a$ then
fixes $q$. If such an $a$ acts trivially on all of $\Omega$,
apply it to $\lambda^{1/n}\in\Omega$ for every $n$.
Their leading terms force $\psi(a/n)=1$, so
$a\in\bigcap_n n\widehat\Z=0$. This proves faithfulness.

Let $\rho_0(g)$ act coefficientwise on $\mathcal P$, fixing
every $q^{1/n}$. It restricts to the continuous tangential
action on $\Omega$ supplied by
$\operatorname{Spec}\Q((q))\to U$, $\lambda\mapsto\lambda(q)$.
The parameter $q$ is a rational formal uniformizer; it need
not belong to $\Omega$. Directly on each Puiseux monomial,
\[
 g\bigl(\psi(a/n)\bigr)=\psi\bigl(\chi(g)a/n\bigr)
\]
proves the displayed conjugation identity. Both types of
operator preserve the selected valuation, hence its completed
cusp $c$.
\end{proof}

\begin{claim}{D.11}{Proved}
There is an abstract semilinear action $\rho$ of $G_\Q$ on
$\mathcal U$ that normalizes no deck group of a finite
hyperbolic quotient. Its $G_\Q^\circ$-fixed set on the
completed pro-cover is precisely the single cusp $c$ of D.10.
This answers the second demand of D5. It does not settle
D5's universal fixed-point assertion.
\end{claim}
\begin{proof}
Fix an odd prime $\ell$ and distinct rational primes
$p_1,p_2,\ldots$ different from $\ell$. Compatible radicals
give Kummer cocycles and their continuous product
\[
 \kappa:G_\Q\longrightarrow V:=\prod_{i\ge1}\Z_\ell,
 \qquad
 \kappa(gh)=\kappa(g)+\chi_\ell(g)\kappa(h).
\]
We have chosen a compatible basis of the Tate twist in
writing these coordinates. Its restriction to $G_\Q^\circ$
is a continuous surjective homomorphism. Indeed, for finitely
many $p_i$ and $m\ge1$, put
\[
 K=\Q(\mu_{\ell^m}),\qquad
 L=K(p_i^{1/\ell^m}).
\]
The extension $K/\Q$ is unramified at the $p_i$.
Valuations there show that their classes in
$K^\times/(K^\times)^{\ell^m}$ are independent of order
$\ell^m$. Kummer theory gives
$T:=\operatorname{Gal}(L/K)\simeq(\Z/\ell^m)^r$.
The extension $L/\Q$ is Galois, and a lift of a cyclotomic
automorphism with character $2$ modulo $\ell$ acts on $T$
by that scalar. Its commutators with $T$ exhaust $T$, since
$\chi_\ell-1$ is a unit. As the quotient by $T$ is abelian,
$T$ is exactly the commutator subgroup. Hence
$L\cap\Q^{\mathrm{ab}}=K$, which proves the required
surjection from $G_\Q^\circ$ at every finite level.
Compactness gives the surjection onto $V$.

Let $e_i\in V$ be the coordinate vectors. They are linearly
independent in $V\otimes_{\Z_\ell}\Q_\ell$. Extend them to
a vector-space basis and choose a linear functional taking
$e_i$ to $\ell^{-i}$. Its restriction is a
$\Z_\ell$-linear map
\[
             f:V\longrightarrow\Q_\ell,
                   \qquad f(e_i)=\ell^{-i}.
\]
This basis choice uses the axiom of choice and specifies an
abstract cocycle; it is not an effective algorithm for
evaluating the cocycle on arbitrary Galois elements.
It is onto, because its image contains
$\ell^{-i}\Z_\ell$ for every $i$. Its kernel is dense:
realize any finitely many prescribed coordinates by a
finitely supported $v$. Choose $j$ outside those coordinates
so large that $\ell^j f(v)\in\Z_\ell$, and use
$v-\ell^j f(v)e_j$. This has the prescribed coordinates
and is killed by $f$.

Regard $\Q_\ell$ as the subgroup of $\mathbb A_f$ supported
at $\ell$ and set $a_g=f(\kappa(g))$. Linearity gives
$a_{gh}=a_g+\chi(g)a_h$. Therefore
\[
                   \rho(g)=T_{a_g}\rho_0(g)
\]
is an action by D.10, semilinear for the specified usual
action on $\overline\Q$. Moreover
\[
 a(G_\Q^\circ)=\Q_\ell,
 \qquad K_0:=\ker(a|_{G_\Q^\circ})
                     \text{ is dense in }G_\Q^\circ.
\]
For density, a continuous surjective homomorphism of compact
groups is an open quotient map; apply this to $\kappa$ and
the dense kernel of $f$.

The orbit of $\lambda$ is infinite: $\rho_0$ fixes this
rational coordinate and D.10 gives
\[
 \{\rho(g)\lambda:g\in G_\Q^\circ\}
       \simeq \Q_\ell/\Z_\ell.
\]
We show why this rules out \emph{every} normalized finite
hyperbolic quotient, without using a congruence-quotient
comparison. If a deck group $H$ were normalized, the action
would descend to its hyperbolic quotient $W$ and the field
$F=\Omega^H$. Choose a number field $K$ defining $W$ and
finitely many function-field generators. For $g\in G_K$,
compare the descended semilinear automorphism to the
coefficient action of that $K$-model. Their quotient is in
the finite group $\operatorname{Aut}_{\overline\Q}(W)$.
Thus those generators have finite $G_K$-orbits, and then
finite $G_\Q$-orbits. Constants also have finite orbits, so
every element of $F$ does. Each element of $\Omega$ is
algebraic over $F$. Its orbit is among the roots of the
finitely many conjugates of a polynomial equation over
$F$, hence is finite. This contradicts the orbit of
$\lambda$. No continuity was presumed in this argument.

It remains to determine the fixed set. Every $T_a$ and
$\rho_0(g)$ fixes $c$, so $c$ is fixed by $\rho$.
If a completed geometric closed pro-point $P$ is fixed by
$\rho(G_\Q^\circ)$, it is fixed by $\rho_0(K_0)$.
Its stabilizer under $\rho_0$ is closed: take cofinal
characteristic finite levels, descending over $\Q$ under
$\rho_0$. At each level an algebraic closed point has an
open and closed Galois stabilizer, and the stabilizer of
$P$ is their intersection. Density of $K_0$ implies that
$P$ is fixed by $\rho_0(G_\Q^\circ)$.

The same Kummer argument as in D.7 proves that the latter
fixed set is $\{c\}$. Explicitly, if $P\ne c$, choose a
characteristic covering level separating their projections
and dominating the degree-nine Fermat cover
$x^3+y^3=1$, $\lambda=x^3$, of $U$. Its smooth proper
compactification has positive genus. More explicitly, if
$N\subset\Pi_U$ is this characteristic subgroup, then
$N\rho_0(G_\Q)\subset\pi_1(U_\Q)$ defines a finite
\'etale curve over $\Q$; take its smooth proper
compactification over $\Q$. The two projected
$\Q^{\mathrm{ab}}$-points have the same proper section
class, since the common action fixes both pro-points.
Their Abel--Jacobi difference has zero Kummer class and
therefore vanishes by F.3, based on the precise abelian
Kummer theorem of \citet[Theorem~3.1]{F_Tsu2020}.
Distinct points on a positive-genus smooth proper curve
cannot have principal difference: that would give a
degree-one map to $\mathbb P^1$. The contradiction proves
$P=c$.
\end{proof}

The constructed action lies outside the finite-quotient
normalization cases covered by R19, while satisfying the
predicted fixed-point conclusion itself. Thus automatic
finite-quotient descent cannot reduce all abstract actions
to B. This construction alone neither proves nor refutes the universal
assertion in B$'$, and does not exclude every logically different
reduction to B. The separate construction in D.12--D.14 below refutes
the universal assertion by an action with empty fixed set.


% Root's ultralimit construction, independently checked and written by the Part F agent.
\subsection{A discontinuous action with no cyclotomic fixed pro-point}

\begin{claim}{D.12}{Proved}
For $U=\mathbb P^1_\Q\setminus\{0,1,\infty\}$, put
$E=\pi_1(U)$ and $\Pi=\pi_1(U_{\overline\Q})$, and fix a
decomposition embedding $G_{\Q_3}\hookrightarrow G_\Q$.
There is an abstract, provably discontinuous section
\[
                         s:G_\Q\longrightarrow E
\]
whose restriction to $G_{\Q_3}$ is geometrically conjugate
to the continuous point section at the explicit transcendental
point
\[
                    a=1+\sum_{r=1}^{\infty}3^{r!}
                       \in U(\Q_3).
\]
It induces an abstract semilinear action on the universal
pro-cover and its completed proper levels, normalizing the
original geometric deck group $\Pi$.
\end{claim}
\begin{proof}
First verify transcendence, rather than assume a generic local
point. Put $A_N=1+\sum_{r=1}^N3^{r!}$. Then
\[
 v_3(a-A_N)=(N+1)!,\qquad
                   1<A_N<2\cdot3^{N!}.
\]
If a nonzero $P\in\Z[T]$ of degree $d$ vanished at $a$,
then $3^{(N+1)!}$ would divide $P(A_N)$. The latter integer
is nonzero for all sufficiently large $N$, since the $A_N$
are distinct. On the other hand
$|P(A_N)|_\infty\leq C(2\cdot3^{N!})^d$ for a fixed
$C$. This contradicts $(N+1)!>dN!+O(1)$. Thus $a$ is
transcendental over $\Q$. It is a $3$-adic unit, and
$v_3(a-1)=1$, so it is an interior point of $U$.

Construct positive integers $x_n\geq2$ recursively. Choose a
prime $q_n\ne3$ which divides none of $x_1,\ldots,x_{n-1}$
and impose
\[
 \begin{split}
 x_n&\equiv a\pmod{3^n},\\
 x_n&\equiv q_n\pmod{q_n^2},\\
 x_n&\equiv1\pmod{q_i}\quad(1\leq i<n).
 \end{split}
\]
The primes $q_n$ are necessarily distinct, and the Chinese
remainder theorem gives such an integer. One may always choose
the least permitted prime and the least positive CRT representative.
That representative is neither zero nor one because of its
congruence modulo $q_n^2$. The first congruence is effective:
only terms $3^{r!}$ with $r!<n$ contribute. We have
\[
           x_n\longrightarrow a\text{ in }\Q_3,
          \qquad v_{q_i}(x_j)=\delta_{ij}\quad(i,j\geq1).
\]
For $j<i$ the valuation assertion is the choice of $q_i$;
for $j>i$ it is the third congruence. In particular the
square classes of the $x_n$ are independent. The exact
twelve-term CRT and valuation certificate is in
\texttt{code/F\_D5\_ultralimit\_output.json}; its role is
only to check this finite arithmetic input, not the
ultrafilter or the fixed-point argument.

Choose any continuous representatives $s_n$ of the global
point sections at $x_n$. Fix a nonprincipal ultrafilter
$\mathfrak u$ on $\mathbb N$ and, for every $g\in G_\Q$,
define
\[
                         s(g)=\lim_{n\to\mathfrak u}s_n(g)
                              \quad\hbox{in }E.
\]
Compact Hausdorffness gives this limit. Continuity of
multiplication in $E$ proves
$s(gh)=s(g)s(h)$; continuity of $E\to G_\Q$ gives
$\operatorname{pr}(s(g))=g$. Thus $s$ is an abstract section.
This argument does not assert continuity of a limit of
continuous maps with the pointwise topology.

We prove exactly the local compatibility required below.
Characteristic-zero geometric base change identifies
$\pi_1(U_{\Q_3})$ with $E\times_{G_\Q}G_{\Q_3}$,
so the local point section $s_a$ is viewed as a homomorphism
$G_{\Q_3}\to E$.
For a finite arithmetic quotient $\varphi:E\twoheadrightarrow F$,
let $P=\varphi(\Pi)$. Its finite \emph{\'etale} Galois
cover $W\to U$ has constant quotient
$W/P=U\times\Spec K_0$, where $K_0/\Q$ is the finite
constant extension with group $F/P$. This is the ordinary
arithmetic fundamental-group covering correspondence
\citep[Expos\'e~IX, \S6]{E_SGA1}.
For $b$ sufficiently close to $a$ in $U(\Q_3)$, Hensel's
lemma gives a Galois-equivariant, $F$-equivariant bijection
\[
                         W_b(\overline\Q_3)
                              \xrightarrow{\sim}
                         W_a(\overline\Q_3)
\]
inducing the identity on the indicated constant quotient.
Here is the uniformity in that statement. Choose a finite
Galois extension $L/\Q_3$ splitting the finite fibre $W_a$.
The \emph{\'etale} inverse-function form of Hensel's lemma
gives, through every point of that fibre, a unique local
$L$-analytic sheet over a sufficiently small ball about $a$.
There are finitely many sheets, deck transformations, and
elements of $\Gal(L/\Q_3)$, so one common ball works.
Uniqueness of these sheets makes all their actions compatible.
The constant coordinates remain their original algebraic
values. This proves the asserted equivariance and the identity
on $W/P$.

In torsor coordinates, such a bijection is multiplication by
an element of $P$, since it is the identity modulo $P$.
Consequently for all sufficiently large $n$ the restrictions
$\varphi s_n|_{G_{\Q_3}}$ and
$\varphi s_a|_{G_{\Q_3}}$ are conjugate by $P$.
This is geometric conjugacy, not merely conjugacy by an
arbitrary arithmetic element of $F$. The $P$-conjugacy orbit
of the latter homomorphism is a finite closed subset of
$F^{G_{\Q_3}}$. Its ultralimit therefore still belongs to
that orbit. For each $\varphi$ there is a geometric
conjugator at this finite level.

The set of elements $\gamma\in\Pi$ satisfying
\[
 \varphi\big(\gamma s(g)\gamma^{-1}\big)=\varphi(s_a(g))
                    \qquad(g\in G_{\Q_3})
\]
is nonempty and closed in $\Pi$. Any finitely many such
conditions have a common refinement given by a finite
quotient of $E$. Compactness of $\Pi$ supplies one
$\gamma$ satisfying all of them. Finite quotients separate
elements of $E$, so
\[
               \gamma s|_{G_{\Q_3}}\gamma^{-1}=s_a
\]
as actual homomorphisms. In particular the local restriction
of the abstract section is continuous.

Its global discontinuity is independently visible in one
finite quotient. The cover $z^2=t$ gives a continuous character
$\tau:E\to\F_2$ such that
$\tau s_n=\kappa_2(x_n)$. Independence of the square classes
gives a continuous surjection
\[
 \kappa:G_\Q\twoheadrightarrow\prod_{n\geq1}\F_2,
                 \qquad g\longmapsto(\kappa_2(x_n)(g))_n.
\]
Indeed it is surjective on every finite product by Kummer
theory, and compactness gives surjectivity onto the product.
The ultrafilter functional
\[
 \epsilon:\prod_n\F_2\longrightarrow\F_2,
              \qquad (v_n)\longmapsto\lim_{n\to\mathfrak u}v_n
\]
is additive. It kills the dense direct sum but sends the
constant-one vector to one, so it is discontinuous. A
continuous surjection of compact Hausdorff groups is a
quotient map; hence $\epsilon\kappa$ is discontinuous.
By construction $\tau s=\epsilon\kappa$, proving the
claimed discontinuity without any section-conjecture input.

Realize $E$ as the arithmetic automorphism group of the
universal pro-cover of $U$. Each $s(g)$ is a semilinear
automorphism, and the homomorphism law gives the required
abstract action. It normalizes $\Pi$, which is normal in
$E$. At every finite covering field the automorphism extends
uniquely to its proper normalization; compatibility gives
the action on the completed pro-cover. The stabilizer of the
function $\sqrt t$ is the nonopen subgroup
$\ker(\epsilon\kappa)$, so the induced action itself is
discontinuous. The explicit sequence and all finite
approximations are effective; choosing a nonprincipal
ultrafilter is a choice-principle step and does not give an
algorithm for the abstract section.

In the crossed-cocycle form requested in R22, fix a continuous
reference section $\rho_0$ at $1/2$ and put
\[
 a_g=s(g)\rho_0(g)^{-1}
     =\lim_{n\to\mathfrak u}s_n(g)\rho_0(g)^{-1}\in\Pi.
\]
Then
\[
 a_{gh}=a_g\,\rho_0(g)a_h\rho_0(g)^{-1},
                 \qquad a_g\rho_0(g)=s(g).
\]
Thus this explicitly specifies an abstract crossed cocycle
for the reference action, and its twisted action is exactly
the action constructed above. D.14 computes its fixed set
as empty.
\end{proof}

\begin{claim}{D.13}{Proved}
Let $L/\Q_p$ be finite, let
$\chi:G_L\to\widehat\Z^\times$ be the full cyclotomic
character, and put $H=\ker\chi$. If $b\in L^*$ satisfies
\[
            \kappa_p(b)|_H=0
                    \quad\hbox{in }H^1(H,\Z_p(1)),
\]
then $b$ is a root of unity. In particular, if the restrictions
of the point sections of $\G_m/L$ at $b_1,b_2\in L^*$
are geometrically conjugate on $H$, then $b_1/b_2$ is a
root of unity.
\end{claim}
\begin{proof}
Choose compatible $p$-power roots and let
$c:G_L\to\Z_p(1)$ be the resulting continuous Kummer
cocycle. Since $H$ acts trivially on the coefficient module,
every $H$-coboundary is zero. The hypothesis therefore says
$c|_H=0$ literally. The cocycle descends to a continuous
cocycle on the compact abelian quotient
$A=\chi(G_L)$: its values are constant on $H$-cosets,
and the compact quotient topology gives continuity. This
proves the needed inflation assertion directly for the
compact coefficient module.

There is $a_0\in A$ with $\chi_p(a_0)\ne1$, since a finite
extension of $\Q_p$ does not contain all $p$-power roots
of unity. The cocycle identity, evaluated on $a_0a=aa_0$,
gives
\[
 (\chi_p(a_0)-1)c(a)=(\chi_p(a)-1)c(a_0)
                                         \qquad(a\in A).
\]
The right side is a coboundary. Thus
$(\chi_p(a_0)-1)[c]=0$ in $H^1(A,\Z_p(1))$.
Writing the nonzero scalar as $p^m$ times a $\Z_p$-unit
shows that $p^m\kappa_p(b)=0$ after inflation to $G_L$.
This is the elementary cocycle calculation often called
Sah's lemma; no vanishing assertion for a nonabelian
quotient is being used.

It follows at every finite Kummer level that
$b^{p^m}\in (L^*)^{p^n}$ for every $n$. We record the local
intersection precisely. Write
\[
 L^*=\varpi^{\Z}\times\mu_{q-1}(L)\times U^1,
               \qquad U^1=1+\mathfrak m_L,
\]
where $q$ is the residue-field cardinality. Hensel's lemma
gives the displayed prime-to-$p$ roots of unity. The group
$U^1$ is pro-$p$, since its quotients modulo
$1+\mathfrak m_L^r$ have $p$-power order. An element
divisible by every $p^n$ in $L^*$ has valuation zero;
after removing its unique $\mu_{q-1}$ part, its image in
each finite $p$-group quotient of $U^1$ is a $p^n$-th
power for every $n$, and therefore identity. The filtration
separates points. Hence
\[
                    \bigcap_{n\geq1}(L^*)^{p^n}
                             =\mu_{q-1}(L).
\]
Thus $b^{p^m}$ is a root of unity, and so is $b$.

The final assertion is the Kummer description of
$\pi_1(\G_m/L)=\widehat\Z(1)\rtimes G_L$ with the
splitting at $1$. Point-section classes are Kummer classes
and their difference is the class of $b_1/b_2$
\citep[\S2.2 and \S8.1--8.3, especially Proposition~26]{F_Stix2008}.
Restrict to $H$ and apply the result just proved.
\end{proof}

\begin{claim}{D.14}{Refuted}
The universal assertion D5, that every abstract semilinear
$G_\Q$-action on the universal pro-cover has exactly one
$G_\Q^\circ$-fixed completed geometric closed pro-point,
is false. The action of D.12 has \emph{no} such fixed
pro-point, even though it normalizes the original geometric
deck group.
\end{claim}
\begin{proof}
Here $G_\Q^\circ=\ker\chi=G_{\Q(\mu_\infty)}$.
The completed pro-points are exactly the compatible systems
of geometric closed points on the proper finite levels,
as in D.7; the generic scheme point is not part of this
set. Suppose that a completed pro-point $\xi$ were fixed
by $s(G_\Q^\circ)$, and let
$x\in\mathbb P^1(\overline\Q)$ be its image at the
base compactification.

Use the fixed embedding $\overline\Q\hookrightarrow
\overline\Q_3$ underlying the decomposition group in D.12,
and choose a finite $L/\Q_3$ containing the coordinates
of $x$. Put
\[
                 H=G_L\cap G_\Q^\circ
                      =\ker(\chi:G_L\to\widehat\Z^\times).
\]
Conjugate $s$ by the geometric element supplied by D.12,
and conjugate $\xi$ with it. Its base point remains $x$,
and we now have $s|_{G_{\Q_3}}=s_a$.
Consequently $s_a(H)$ fixes a completed pro-point over $x$.

Choose the following unit of the open tripod:
\[
\begin{array}{c|c|c}
 x&f(t)&f(x)\\\hline
 x\in U(\overline\Q)&t&x\\
 0&1-t&1\\
 1&t&1\\
 \infty&t/(t-1)&1
\end{array}
\]
In every case $f$ is regular and nonzero at $x$.
The morphism $f:U\to\G_m$ therefore sends the decomposition
group of the fixed pro-point over $x$ into the point
decomposition group at $f(x)$. At a cusp its geometric
inertia is killed, since the order of $f$ there is zero.
Equivalently, the covers $w^{3^n}=f(t)$ extend
\emph{\'etale} across $x$ and their fibre at $x$ consists
of the $3^n$-th roots of $f(x)$. Evaluating the section
in these covers gives
\[
       \kappa_3\bigl(f(a)/f(x)\bigr)|_H=0
                         \quad\hbox{in }H^1(H,\Z_3(1)).
\]
This is the usual functorial evaluation of units on
sections \citep[Proposition~26 and \S8.3]{F_Stix2008}.
Changing the selected roots changes only a coboundary;
on $H$ that coboundary is zero because the cyclotomic
action is trivial. The argument applies to interior
points and to each cusp, with no supposition that the
original abstract global section is continuous.

Both $f(a)$ and $f(x)$ belong to $L^*$, and D.13 gives
$f(a)/f(x)\in\mu(L)$. The number $f(x)$ is algebraic
over $\Q$, as is every root of unity. Hence $f(a)$
would be algebraic over $\Q$. Each chosen $f$ is an
invertible fractional-linear function over $\Q$, so
this would make $a$ algebraic, contradicting D.12.
Thus the fixed set is empty.
It is the existence part of the asserted unique fixed point
that fails in this example; no assertion of two distinct
fixed points is made.

D.11 and this counterexample address the two demands
of D5 separately: D.11 constructs an action normalizing
no finite-quotient deck group and having exactly one fixed
cusp, while the present action disproves the universal
fixed-point claim. The latter uses discontinuous abstract
actions and proves no counterexample to the continuous
section conjecture in D1--D4.
\end{proof}


The former obligation D.U5 is superseded by D.12--D.14:
the universal fixed-point assertion is false, with an empty
fixed set already possible for an action normalizing the
original deck group. D.11 separately supplies the requested
action normalizing no finite-quotient deck group.

\subsection{The local completion argument: setup and the first stage}
\label{sec:D-local-first-stage}

The next two statements reconstruct Sections 2--4 of the supplied
manuscript \citet{D_OpenAI2026}. They are a checked derivation of that
source, with the comparison maps and the treatment of every inverse
annulus made explicit. No priority is claimed. The conclusion of this
subsection is $L_\ell\subset H$; the passage from this field to the
whole separable closure requires the subsequent second stage.

\paragraph{Fixed fields and the valued point.}
Let $p$ be a prime, let $k/\Q_p$ be finite, and let $X/k$ be smooth,
proper, geometrically connected, and of genus at least two. Set
$K=k(X)$ and fix a separable closure $K^{\mathrm{sep}}$ containing an
algebraic closure $\bar k$ of $k$. Fix a real-valued nonarchimedean
field norm on $K$ extending the norm on $k$, normalized by
$|p|=p^{-1}$, and fix one prolongation to $K^{\mathrm{sep}}$.
Thus its support on $K$ is zero. All fields below have these specified
embeddings and norms. Put
\[
 K_0=\bar k(X),\qquad
 \widetilde K=\operatorname{compositum}_{V\to X}
                    k(V)\subset K^{\mathrm{sep}},
 \qquad E=\widehat{K^{\mathrm{sep}}},\qquad
 H=\widehat{\widetilde K}\subset E.
\]
In the compositum, $V$ runs through connected finite \'etale covers and all their
$K$-embeddings in the fixed separable closure. The displayed inclusion
of $H$ identifies it with the closure of $\widetilde K$ in $E$.
Finite constant-field covers show that $\bar k\subset\widetilde K$.
Consequently
\[
                     \Omega=\widehat{\bar k}\subset H
\]
is a specified complete algebraically closed subfield. Here we use
the usual fact that the completion of an algebraically closed
rank-one valued field is algebraically closed. For
$v=-\log_p|\cdot|$ one has $v(\Omega^\times)=\Q$: this holds on
$\bar k$, and a nonzero limit has the same norm as every sufficiently
close approximant, so completion adds no value.

The inclusion $K_0\hookrightarrow H$ defines an $H$-valued point of
$X_{\bar k}$ supported at its generic point. Together with the fixed
inclusion $\Omega\hookrightarrow H$, it defines a specified
$H$-valued point of $X_\Omega$. Denote its analytic image by
$\xi\in X_\Omega^{\rm an}$. We retain the full valued point,
including its isometric map of completed residue fields
$\mathcal H(\xi)\hookrightarrow H$. Its support on $X_\Omega$ can
be closed, even though its projection to $X_{\bar k}$ is generic;
in particular, $\xi$ can be an $\Omega$-rational point. It cannot
be any of the algebraic points in $X(\bar k)$.

For a finite extension $F/K_0$, normalization gives a finite map of
smooth proper curves over $\bar k$. Its \emph{horizontal indices}
mean the ramification indices at their closed points. These discrete
valuations are trivial on $\bar k$ and have residue characteristic
zero. Fix an odd prime $\ell\ne p$ and define, inside
$K^{\mathrm{sep}}$,
\[
 L_\ell=\operatorname{compositum}\{F: F/K_0\text{ finite Galois,
                       every horizontal index divides }\ell\}.
\]
Then $\widetilde K\subset L_\ell$. Indeed, after adjoining constants
the universal finite \'etale field is still $\widetilde K$: every
finite \'etale cover of $X_{\bar k}$ descends to a finite extension
of $k$, and composition with that constant-field cover gives a finite
\'etale cover of $X$.

\begin{claim}{D.15}{Proved; analytic patching, evaluation, and horizontal indices}
\label{clm:D15}
The following assertions hold in the preceding setup.
\begin{enumerate}[label=(\roman*)]
\item Let $f:C\to X_{\bar k}$ be finite and surjective, with $C$
smooth, proper, connected, of degree $d$ and branch set $B$.
Choose a rational local parameter $t_b$ at each $b\in B$.
There are disjoint open coordinate disks $B_b$ about these points,
avoiding $\xi$, such that the entire inverse disk is a disjoint
union of disks with equations $t_b=z_c^{e_c}$, where
$e_c=e(c/b)$. Suppose $g:Y\to X_\Omega$ is finite and surjective,
with $Y$ smooth, proper and connected, and that there are integers
$m_b$ divisible by all $e_c$ over $b$ and radii
$0<r_b<R_b$ in $|\Omega^\times|$ strictly inside $B_b$ for which
$g^*t_b$ has an analytic $m_b$th root on all of
$g^{-1}A_b$, where $A_b=\{r_b<|t_b|<R_b\}$.
Set
\[
 U=X_\Omega^{\rm an}\setminus
            \bigcup_{b\in B}\{|t_b|\le r_b\}.
\]
There is a connected finite \'etale surjective cover $V_\Omega\to Y$
and an analytic map
\[
     V_\Omega^{\rm an}\big|_U\longrightarrow C_\Omega^{\rm an}
                         \quad\hbox{over }X_\Omega^{\rm an}.
\]
The notation on the left is the inverse image of $U$ under the
composite map to $X_\Omega$.
\item Suppose that this composite $V_\Omega\to X_\Omega$ is
isomorphic over $X_\Omega$ to the base change of a finite map
$V\to X_{\bar k}$ from a smooth proper connected curve. If
$\bar k(V)$ has a $K_0$-embedding into $H$, then so does
$\bar k(C)$. The analytic exterior map need not descend to $\bar k$.
\item For every positive integer $n$, the condition that horizontal
indices divide $n$ is preserved by subextensions, finite composita,
Galois closures, normalized pullback along a finite map of smooth
proper curves, and composition with a finite \'etale upper cover.
In particular, a finite curve cover with horizontal indices dividing
$\ell$ has a $K_0$-embedding of its function field into $L_\ell$.
\end{enumerate}
\end{claim}
\begin{proof}
These are the assertions of Lemmas 3.1, 3.3, 3.4 and Proposition 3.2
of \citet{D_OpenAI2026}; we give the construction and the precise
uses of analytic comparison.

\paragraph{Local normalization, including indices divisible by $p$.}
At a point $c$ above $b$, choose an analytic parameter $u$ and write
$f^*t_b=u^{e_c}a(u)$, with $a(0)\ne0$. Choose an $e_c$th root of
$a(0)$ in $\Omega$. On a sufficiently small disk the inequality
\[
 v\bigl(a(u)/a(0)-1\bigr)>v_p(e_c)+\frac1{p-1}
\]
holds uniformly. The convergent expression
\[
 a(u)^{1/e_c}=a(0)^{1/e_c}
       \exp\bigl(e_c^{-1}\log(a(u)/a(0))\bigr)
\]
therefore supplies the desired unit root. The function
$z_c=u a(u)^{1/e_c}$ has nonzero derivative at $c$, so analytic
inversion gives a smaller disk on which it is a coordinate and
$f^*t_b=z_c^{e_c}$. No condition $p\nmid e_c$ was used.

Do this in disjoint disks at the finitely many points above $b$.
Their complement in the proper analytic curve $C_\Omega^{\rm an}$
is compact. Its image is closed and does not contain $b$.
A small enough base disk therefore has its \emph{entire} inverse
image in those chosen disks. Shrink the base disk once more so that
within each coordinate chart the inverse image is exactly the disk
specified by the bound on $|z_c|^{e_c}$. Shrinking disks form a
neighborhood basis at $b$. Hausdorffness, the finiteness of $B$, and
$\xi\notin X(\bar k)$ let us choose the $B_b$ disjoint from one
another and from $\xi$, with radii in $|\Omega^\times|$.

\paragraph{Patching every sheet.}
\label{clm:D15-patching}
The open sets
\[
 Y_U=g^{-1}U,
 \qquad Y_b=g^{-1}\{|t_b|<R_b\}
\]
cover $Y^{\rm an}$; different $Y_b$ are disjoint and
$Y_U\cap Y_b=g^{-1}A_b$. On $Y_U$ take the pullback of
$C_\Omega^{\rm an}$, a finite \'etale cover of degree $d$ because
all branch points have been removed. On each $Y_b$ take the trivial
cover of degree $d$. Write $q_b^{m_b}=g^*t_b$ on $g^{-1}A_b$.
For every point $c$ above $b$, the corresponding algebra on this
overlap splits as
\[
 Z^{e_c}-g^*t_b
   =\prod_{\zeta\in\mu_{e_c}}
              \bigl(Z-\zeta q_b^{m_b/e_c}\bigr).
\]
The function $q_b$ is a unit there. For distinct roots of unity,
$(\zeta-\zeta')q_b^{m_b/e_c}$ is an invertible analytic function.
This remains true when $p\mid e_c$: its constant factor is nonzero
in the field $\Omega$, independently of its absolute value.
The Chinese remainder theorem thus trivializes every local factor.
Since $\sum_{c\mid b}e_c=d$, it trivializes all $d$ sheets on the
\emph{whole} overlap.

Choose these trivializations and glue the finite locally free
analytic algebras. There are no overlaps between distinct $Y_b$,
so no further triple-overlap compatibility is needed. The resulting
coherent algebra is finite \'etale locally on $Y^{\rm an}$.
For the proper $\Omega$-scheme $Y$, the coherent-sheaf equivalence
of \citet[Appendix A, Theorem A.1(ii)]{D_Poineau2010} algebraizes
this sheaf. Full faithfulness algebraizes its unit and
multiplication and preserves the algebra identities. Its relative
spectrum is finite over $Y$, and it is \'etale because
analytification detects that property
\citep[Fact 5.1.5]{D_Temkin2010}. Choose a nonempty connected
component $V_\Omega$. Its image in connected $Y$ is open and
closed, or equivalently its finite \'etale degree is a positive
locally constant function; hence it is surjective. It is smooth
and proper. The original projection from the pullback over $Y_U$
gives the claimed analytic exterior map. If $B=\varnothing$,
the pullback is already finite \'etale and only this connected
component step is needed.

\paragraph{Evaluation at the specified $H$-point.}
\label{clm:D15-evaluation}
Let $j:\bar k(V)\hookrightarrow H$ be the embedding in (ii).
It gives a generic $H$-valued point of $V$ over the fixed generic
point of $X_{\bar k}$. Together with $\Omega\hookrightarrow H$
it gives an $H$-point of $V\times_{\bar k}\Omega$, and the chosen
isomorphism over $X_\Omega$ transports it to $V_\Omega$.
Its image in $X_\Omega(H)$ is exactly the point defined by
$K_0\hookrightarrow H$ and $\Omega\hookrightarrow H$.
Its analytic image is consequently above $\xi\in U$.

For a complete nonarchimedean extension $H/\Omega$, algebraic and
analytic $H$-points agree
\citep[Fact 5.1.3]{D_Temkin2010}. More explicitly the lifted point
is an analytic point $v$ equipped with an isometric embedding
$\mathcal H(v)\hookrightarrow H$. If $c'$ is its image under the
exterior analytic map, that map supplies
\[
             \mathcal H(c')\longrightarrow\mathcal H(v)
                                      \longrightarrow H.
\]
Thus it supplies an $H$-point of $C_\Omega$ over the same
\emph{valued} point of $X_\Omega$, not merely over the same
underlying topological point. Projecting to $C$, the support lies
over the generic point of $X_{\bar k}$. Since $C\to X_{\bar k}$
is finite, that support is generic in $C$. We obtain a field
embedding $\bar k(C)\hookrightarrow H$ whose restriction to
$K_0$ is the specified inclusion. This proves (ii), without any
descent assertion for the exterior analytic map or its local
trivializations.

\paragraph{Horizontal ramification under the allowed operations.}
\label{clm:D15-horizontal}
At a horizontal point the completed base field is
$A=\bar k((t))$. A finite extension has residue field $\bar k$
and is $\bar k((u))$ for a uniformizer $u$: successive subtraction
of representatives from the coefficient field $\bar k$ gives the
power-series expansion in its complete valuation ring.
Writing $t=u^e a(u)$ with $a(u)$ a unit, Hensel's lemma extracts
an $e$th root of the unit, because this is a residue-characteristic
zero valuation. A change of parameter therefore gives $t=z^e$.
Inside a fixed algebraic closure the extension is
$A(t^{1/e})$, which is cyclic since $\bar k$ contains all roots
of unity. If $e\mid n$, this field lies in $A(t^{1/n})$.

Every completed local factor of a subextension, finite compositum,
or normal closure consequently remains inside that same cyclic
extension. This proves the index assertion for those operations.
For a base change write its local equation as $t=s^r a(s)$,
$r\ge1$. Extracting an $e$th root of the unit shows that each
normalized local factor has ramification index
$e/\gcd(e,r)$, which divides $e$. Finally indices multiply in a
tower and an \'etale upper cover contributes index one. Taking
the Galois closure of a cover whose indices divide $\ell$ and
embedding it in $K^{\mathrm{sep}}$ gives the final assertion.
This proves (iii). It concerns horizontal valuations, and does not
assert that composing two ramified covers preserves the bound.
\end{proof}

\begin{claim}{D.16}{Proved; roots on every inverse collar and $L_\ell\subset H$}
\label{clm:D16}
Let $\ell\ne p$ be the odd prime fixed above.
\begin{enumerate}[label=(\roman*)]
\item For a smooth proper connected curve $C/\Omega$ of genus at
least two, an annular coordinate $S$ on a node annulus of its stable
skeleton acquires an analytic $\ell$th root on the \emph{entire}
inverse annulus in a connected finite \'etale cover of $C$.
One cover can serve a finite collection of node annuli.
\item Let $C/\bar k$ have the same properness, smoothness and genus
hypotheses, and let $b_i$ be finitely many distinct points, with
rational local parameters $t_i$. Inside any sufficiently small
pairwise disjoint coordinate disks $B_i\subset C_\Omega^{\rm an}$
about them, there are concentric collars
\[
 A'_i=\{r_i<|t_i|<R_i\},\qquad
             r_i,R_i\in|\Omega^\times|,\quad 0<r_i<R_i,
\]
and a connected finite \'etale cover $Y\to C_\Omega$ on which
$t_i$ has an analytic $\ell$th root on all of the inverse image
of $A'_i$, for every $i$.
\item With the specified subfields of $E$ in the setup,
\[
                              L_\ell\subset H.
\]
\end{enumerate}
\end{claim}
\begin{proof}
This reconstructs Lemmas 4.1--4.3 and Proposition 4.4 of
\citet{D_OpenAI2026}. We first give the elementary analytic input,
then construct the covers, and finally apply D.15.

\paragraph{Units on complete open disks and annuli.}
An analytic unit on an open disk has constant norm in
$|\Omega^\times|$. On an open annulus $A=\{r<|S|<R\}$ every unit
has a factorization
\[
 u=cS^a(1+h),\qquad c\in\Omega^\times,\quad a\in\Z,
       \quad h\in\mathcal O(A),\quad |h(x)|<1\quad(x\in A).
 \tag{D.16.1}
\]
Here $cS^a$ is a single term of its Laurent expansion, and the
factorization holds on the entire open annulus.
To verify this, restrict to a circle $|S|=\rho$ with
$\rho\in|\Omega^\times|$. Its Laurent Gauss norm is
multiplicative. The norm of a nonzero analytic function is attained
by one coefficient term and belongs to $|\Omega^\times|$.
Normalize $u$ and $u^{-1}$ to norm one with reciprocal constants.
Their reductions are inverse units in
$\widetilde\Omega[s,s^{-1}]$, whose units are monomials.
Exactly one Laurent term of $u$ therefore dominates at this radius.
The dominant exponent cannot change as the radius moves through
the annulus. Indeed, convergence on any compact subannulus leaves
only finitely many candidate dominant terms. A change forces a tie
of two such terms; their coefficient valuations are rational, so
the tie radius lies in $|\Omega^\times|$, contradicting the
preceding uniqueness. One term dominates on every compact
subannulus, and hence on their union, giving (D.16.1).
The Gauss inequality also gives the strict bound on $h$ away from
the skeleton. For a disk the corresponding reduction ring on
closed subdisks is $\widetilde\Omega[s]$, with only constant
units. The same argument shows constant norm; evaluating at a
rational point puts that norm in $|\Omega^\times|$.

Because $1/\ell\in\Z_p$, the binomial series for
$(1+h)^{1/\ell}$ converges on every compact subannulus: the
continuous norm of $h$ has a maximum strictly below one there.
These roots agree on overlaps because they are given by the same
series. Also $c$ has a root in $\Omega$. Thus if $S$ acquires an
$\ell$th root on an analytic inverse image, every unit $u$ on $A$
does too, including when $a<0$. More generally the unit already
has a root on $A$ whenever $\ell\mid a$.

\paragraph{The exact lifting input.}
A split stable model of a proper curve determines a finite metric
graph, its skeleton, and at each vertex the smooth normalization
$D_v$ of the residue component, with a marked point for each
incident edge end. A loop has two distinct marks. The tube of a
node is the annulus retracting to the corresponding open edge.
For $g(v)=g(D_v)$, stability is
\[
                         2g(v)-2+\operatorname{val}(v)>0.
 \tag{D.16.2}
\]
The residue field $\widetilde\Omega$ is algebraically closed of
characteristic $p$, and all edge lengths lie in $\Q$.
We use \citet[Theorem 7.7]{D_ABBR2015} for this fixed target
curve. Its input is a tame covering of metrized complexes:
a finite harmonic graph map, finite tame residue-curve maps with
ramification indices equal to the edge expansion factors, all
fibres of all marked points represented by edge ends, and zero
augmented graph ramification away from punctures
\citep[Definitions 2.19 and 2.21]{D_ABBR2015}.
We take no punctures, so the lift has empty branch locus and is
finite \'etale; the lifted vertex set and skeleton are exactly the
inverse images of the target ones
\citep[Definitions 4.25 and 4.31]{D_ABBR2015}.
The formulation of Theorem 7.7 includes loop edges.

\paragraph{A nonseparating stable edge.}
Let $e$ be the edge for the given annulus $A$, and first suppose
that $e$ is nonseparating. Choose a simple circuit through $e$.
At a vertex on this circuit let $P,Q$ be its two circuit marks
on $D_v$. Multiplication by $\ell$ on $\operatorname{Pic}^0(D_v)$
is surjective over the algebraically closed residue field.
Thus choose a degree-zero divisor $T$ for which
$P-Q+\ell T$ is principal, say the divisor of $f_v$.
The normalized cover $z^\ell=f_v$ is connected, because its
order at $P$ is congruent to $1$ modulo $\ell$, and it has degree
$\ell$. It is totally ramified at $P,Q$ and unramified elsewhere.
Each other marked fibre therefore has $\ell$ distinct points.
At every vertex outside the circuit take $\ell$ identity copies
of its residue curve.

Over each circuit edge take one edge with expansion factor
$\ell$, and length the target length divided by $\ell$.
Over every other edge take $\ell$ edges of factor one, matching
all marked fibres at both ends. This prescription includes chords
between circuit vertices and both ends of a loop. At each circuit
vertex the total degree in every target direction is $\ell$;
at an identity-copy vertex it is one. The residue formula
\[
       2g(D'_v)-2=\ell(2g(D_v)-2)+2(\ell-1)
 \tag{D.16.3}
\]
accounts for exactly the two ramified circuit directions.
There is consequently zero augmented graph ramification.
All residue ramification is tame, since $\ell\ne p$.
The resulting covering is connected: lift a path from any vertex
outside the circuit until its first meeting with the circuit;
its endpoint is on the connected lifted circuit. The source is
stable as well. At a circuit vertex its valence is
$2+\ell(\operatorname{val}(v)-2)$ and (D.16.3) gives
\[
 2g(v')-2+\operatorname{val}(v')
       =\ell\bigl(2g(v)-2+\operatorname{val}(v)\bigr)>0.
\]
Identity copies preserve (D.16.2).

Lift this tame covering by the stated theorem. We explain why its
inverse over $A$ consists of the \emph{whole} source node annuli,
including all their off-skeleton points. An off-skeleton open ball
maps into the complement of the target skeleton, by the equality
of inverse skeleta. Its connected image lies in one target
off-skeleton ball. Continuity sends its attachment point to that
target ball's attachment point, since the latter is the unique
point of the target skeleton in the ball's closure. Thus the
retractions commute on these balls and on the skeleton.
The inverse of the retraction tube of the open edge $e$ is exactly
the union of the corresponding source open-edge tubes. These are
the full node annuli above $A$.
On each one the Laurent exponent of the pullback of $S$ is
$\ell$ or $-\ell$, by the prescribed expansion factor and the
choice of orientation. Formula (D.16.1) supplies an analytic
$\ell$th root there. The annuli are disjoint, so these functions
give a root on the entire inverse image of $A$.

\paragraph{A separating stable edge, with every lifted edge treated.}
Suppose now that $e$ separates the graph. Each of the two sides,
retaining the boundary half-edge, contains either a circuit or a
positive-genus vertex. Otherwise it would be a tree with $N$
genus-zero vertices and total valence $2N-1$, including that
half-edge. Stability requires total valence at least $3N$, a
contradiction.

Construct a connected degree-$\ell$ unramified covering of each
side. If it contains a circuit, take the cyclic graph cover from
a surjection of its graph fundamental group onto $\Z/\ell\Z$,
with identity residue maps. If instead a positive-genus vertex is
used, choose a line bundle $L$ of exact order $\ell$ on its residue
curve. Such bundles exist because its Jacobian has nonzero
$\ell$-torsion. A trivialization $L^{\otimes\ell}\simeq\mathcal O$
makes $\bigoplus_{j=0}^{\ell-1}L^{\otimes j}$ a cyclic algebra.
Its spectrum is finite \'etale, since $\ell$ is invertible,
and connected: the nontrivial summands have no global sections
as nontrivial degree-zero line bundles, so the algebra has global
sections only the residue field and admits no nontrivial idempotent.
Use this connected cover at the chosen vertex and $\ell$ identity
copies elsewhere, matching all marked fibres. Every copy connects
to the distinguished vertex by lifting a path to its first
occurrence, so this covered side is connected.

Each side has $\ell$ preimages of its boundary mark. Match these
in pairs by $\ell$ factor-one edges. Both covered sides are
connected, so each of the new edges above $e$ is nonseparating:
removing it leaves another edge joining the two sides. The residue
maps are unramified and all edge factors are one. Harmonicity
and zero augmented ramification follow directly, or from the
unramified Riemann--Hurwitz formula. At a vertex of local degree
$d$ one has
\[
 2g(v')-2+\operatorname{val}(v')
                =d\bigl(2g(v)-2+\operatorname{val}(v)\bigr)>0,
\]
so the source triangulation is stable. Lift the covering; the
result is finite \'etale and every annulus over $A$ is now a
nonseparating stable-node annulus.

Apply the preceding circuit construction to each of these
$\ell$ annuli separately. Take a nonempty connected component
of the finite fibre product of the resulting finite \'etale
covers over this intermediate curve. Its projection to each
connected factor is finite \'etale and nonempty, hence
surjective. Every previously constructed root persists by
pullback on its entire prescribed inverse annulus. On each
intermediate factor-one annulus the original $S$ is an analytic
unit; (D.16.1) gives its root once that annulus's coordinate
has a root. This proves (i) for the original separating edge.
Taking a further connected component of a finite fibre product
treats any finite list of original node annuli simultaneously.

\paragraph{Moving arbitrary small collars into stable node annuli.}
All the algebraic data $C,b_i,t_i$ in (ii) descend to a finite
extension of $k$. After another finite extension choose a split
semistable model with the $b_i$ as disjoint smooth sections.
Around the $i$th section let $u_i$ be an integral local parameter.
Successive blowups of its reductions produce node annuli
\[
                  |\pi|^{j+1}<|u_i|<|\pi|^j,
 \tag{D.16.4}
\]
where $\pi$ is a uniformizer of that finite local field.
On a sufficiently small disk $t_i/u_i$ is an analytic unit,
with constant norm by (D.16.1)'s disk assertion. Thus (D.16.4)
is also a concentric annulus for $t_i$. Taking $j$ large puts
it inside the originally prescribed $B_i$. Carry out these
modifications for all the sections and call the resulting
proper semistable model $\mathcal C$. Forget its markings.
The model need not be stable; its newly introduced nodes
are precisely why the following input is required.

We use the model-domination consequence of resolution of
nonsingularities in
\citet[Theorem A and Proposition 2.4(iv)]{D_MochizukiTsujimura2024}.
For a hyperbolic curve over a finite extension of $\Q_p$, the set
of all primes contains $p$ and has the cardinality required there;
its residue field is algebraic over $\F_p$. Applied to the
prescribed compactified model $\mathcal C$, those results give,
after a finite base extension, a connected finite \'etale cover
whose compactified stable model maps to $\mathcal C$.
Pass to split reduction and a geometric connected component.
After base change to $\Omega$ write the curve map as
$f:Y_0\to C_\Omega$, and the compatible model map as
$\mathcal Y_0\to\mathcal C$. We require a model morphism,
not a finite model morphism. The source stable model has only
finitely many nodes.

Fix one target node annulus $A_i\subset B_i$ from (D.16.4).
Choose an irrational real number $a$ strictly between its
additive endpoint valuations. We claim that every point above
the whole level $\{x\in A_i:v(t_i(x))=a\}$ lies in a source
node annulus whose node maps to the target node.
The specialization maps commute with the model morphism.
A source point reducing to a component generic point is the
corresponding model vertex, a type-2 point whose value group is
$|\Omega^\times|$. It gives rational valuation to $f^*t_i$,
so cannot lie above this level. A source point reducing to a
smooth closed point lies in its open disk tube. If any such
point lies above the level, its reduction maps to the target
node; therefore its \emph{entire} disk tube maps into $A_i$.
On that disk $f^*t_i$ is a unit, hence has constant norm in
$|\Omega^\times|$, again excluding the irrational level.
The remaining possible reductions are source nodes mapping
to the target node, proving the claim.

Let $N_i$ be the union of all these source node annuli. It is a
finite union of open sets. The complement
$Y_0^{\rm an}\setminus N_i$ is compact, since $Y_0$ is proper.
Choose a closed concentric subannulus $A_i^0\subset A_i$
containing the chosen level in its interior. On the compact set
\[
 (Y_0^{\rm an}\setminus N_i)\cap f^{-1}(A_i^0)
\]
the continuous real function $v(f^*t_i)$ has compact image
which avoids $a$. Choose an interval about $a$ with rational
endpoints, contained in the valuation interval of $A_i^0$ and
disjoint from that compact image. It defines a collar
$A'_i\subset A_i$ whose \emph{entire} inverse image is contained
in $N_i$. This compact-complement step controls every sheet.

Every source annulus in $N_i$ maps into $A_i$ by specialization
compatibility, so $f^*t_i$ is a unit on that full annulus.
Apply (i) to all source annuli in all the $N_i$ at once.
Their coordinates acquire $\ell$th roots in one connected
finite \'etale cover $Y\to Y_0$. The unit factorization
(D.16.1) then supplies roots of $f^*t_i$ on their entire inverse
images, and hence on all of the inverse image of $A'_i$.
The composite $Y\to C_\Omega$ is finite \'etale. This proves (ii).

\paragraph{The first-stage field inclusion, with the specified copy.}
Let $F/K_0$ be one of the finite Galois extensions defining
$L_\ell$, and let $C\to X_{\bar k}$ be its smooth proper curve
cover. If it is unramified then $F\subset\widetilde K\subset H$.
Otherwise choose the disjoint branch disks of D.15, avoiding
$\xi$. Assertion (ii), applied to $X_{\bar k}$, supplies a
connected finite \'etale $Y\to X_\Omega$ and collars on whose
entire inverse images the branch parameters have $\ell$th roots.
D.15(i), with $m_b=\ell$, yields a connected finite \'etale
$V_\Omega\to Y$ and its analytic exterior map to $C_\Omega$.
The composite $V_\Omega\to X_\Omega$ is finite \'etale.

For a proper connected scheme over an algebraically closed field,
extension to an algebraically closed overfield induces an
equivalence of finite \'etale cover categories
\citep[Expos\'e X, Corollary 1.8]{E_SGA1}. Apply this to
$\bar k\subset\Omega$: the composite descends, up to
isomorphism over $X_\Omega$, to a connected finite \'etale
$V\to X_{\bar k}$. Its function field has a $K_0$-embedding
into $\widetilde K$, by the definition and constant-field
description of that universal \'etale field. D.15(ii) evaluates
the exterior analytic map at this point and gives a
$K_0$-embedding $j:F\hookrightarrow H\subset E$.

It remains to identify its image with the previously specified
subfield $F\subset K^{\mathrm{sep}}\subset E$. Normality is
essential for this last statement: every minimal polynomial
over $K_0$ of an element of $F$ splits in the specified copy
of $F$. Hence $j$ sends every element to an element of that
copy, and $j(F)\subset F$. Equality follows from the equal
finite degrees. Since $j(F)\subset H$, the specified $F$
itself lies in $H$. Taking the compositum over these fields
proves $L_\ell\subset H$ inside $E$, as asserted.
\end{proof}

% Annotated reconstruction of the supplied R24, Sections 5--6.
% Uses the fixed valued fields and first-stage statements of D.15--D.16.
\subsection{Arbitrary collar roots and the specified completion}

The next two claims reconstruct the second stage of the supplied
R24 argument \citep[Sections~5--6]{D_OpenAI2026}.
We retain the fields, chosen valuation and embeddings of
D.15--D.16, including $\Omega\subset H\subset E$, the full
$H$-valued point above $\xi$, and the fixed odd prime
$\ell\ne p$. In particular $v=-\log_p|\cdot|$ and
$L_\ell\subset H$ by D.16. The roots below are analytic over
$\Omega$; no descent of those roots to $\bar k$ is required.

\begin{claim}{D.17}{Proved}\label{clm:D17}
Let $m\ge1$, choose an integer
\[
 D>v_p(m)+\frac1{p-1},
\]
and let $\alpha\in\bar k^*$ satisfy $a=v(\alpha)>2D$.
There are a smooth proper connected curve $Q/\bar k$ and a
finite surjective map $Q\to\mathbb P^1_{\bar k}$ whose
horizontal ramification indices divide $\ell$, such that the
coordinate $T$ has an analytic $m$-th root on the entire inverse
image in $Q_\Omega^{\mathrm{an}}$ of
\[
                   A^\circ=\{D<v(T)<a-D\}.
\]
The conclusion holds for every $m$, including powers of $p$
and the case $p=2$.
\end{claim}
\begin{proof}
\emph{Division on the uniformizing torus, on every sheet.}
We first prove the precise period-lattice statement used below,
corresponding to R24, Lemma~5.1.
Suppose $q:T^{\mathrm{an}}\to J^{\mathrm{an}}$ is a toric
uniformization with free discrete kernel $\Lambda\subset T(\Omega)$,
and $q$ is a topological covering and a local analytic isomorphism.
For an analytic map $a_A:A\to T^{\mathrm{an}}$, form
\[
 B=A\times_{q\circ a_A,J^{\mathrm{an}},[m]}J^{\mathrm{an}},
       \qquad \pi:B\to A,\quad b:B\to J^{\mathrm{an}}.
\]
Every character value $\chi\circ a_A\circ\pi$ has an analytic
$m$-th root on all of $B$.
Indeed, locally lift $b$ to $t:B\to T^{\mathrm{an}}$.
The relation $[m]b=q\circ a_A\circ\pi$ gives
\[
                  t^m/(a_A\circ\pi)\in\Lambda.
\]
This is locally constant because the kernel is discrete.
Replacing $t$ by $t\rho$, $\rho\in\Lambda$, changes this
period by $\rho^m$. Thus its class in $\Lambda/\Lambda^m$
is globally defined and gives a finite clopen decomposition
$B=\coprod_{\bar\lambda}B_{\bar\lambda}$.
Choose a representative $\lambda$ of each nonempty class.
On $B_{\bar\lambda}$, a locally constant period correction
makes every local lift satisfy
\[
                         t^m=\lambda(a_A\circ\pi).
\]
Two corrected lifts on an overlap differ by a period $\rho$
with $\rho^m=1$. The free lattice has no torsion, so $\rho=1$.
The corrected lifts therefore glue to one analytic map
$t_{\bar\lambda}:B_{\bar\lambda}\to T^{\mathrm{an}}$.
Choose $d_\lambda\in\Omega^*$ with
$d_\lambda^m=\chi(\lambda)$. Then
$\chi(t_{\bar\lambda})/d_\lambda$ is the required root.
The subsets are disjoint, so these functions give one function
on all of $B$. Neither local connectedness of a chosen sheet
nor a prime-to-$p$ restriction on $m$ has been inserted.

\emph{The auxiliary curve and its split stable model.}
Let $Z$ be the smooth projective normalization of
\[
                         y^\ell=\frac{T(T-1)}{T-\alpha}.
\]
The four branch orders at $0,\alpha,1,\infty$ are
$1,-1,1,-1$. Thus $Z\to\mathbb P^1$ is a connected cyclic
cover of degree $\ell$, totally ramified at those points and
unramified elsewhere. Riemann--Hurwitz gives
\[
                  2g(Z)-2=-2\ell+4(\ell-1),
                     \qquad g(Z)=\ell-1.
\]
All algebraic data descend to a finite extension of $k$.
Over its valuation ring, blow up the line model along
$(T,\alpha)$. The resulting semistable model has two rational
components and one node with chart $TU=\alpha$; the inner
component has coordinate $S=T/\alpha$.
The branch sections are disjoint and smooth. The Kummer
function reduces to $T-1$ on the outer component and
$-S/(S-1)$ on the inner component. Their degree-$\ell$
normalizations are rational, each with exactly two totally
ramified points.

At the node the same function is
\[
                           \frac{T-1}{1-U},
\]
a unit whose reduction is $-1$. Since $\ell$ is invertible,
normalization is \emph{\'etale} there, with $\ell$ split node
charts after a finite constant extension. Away from the node
and the branch sections the unit criterion gives the same
\emph{\'etale} assertion. At a branch section the function or
its inverse is a simple relative parameter, whose Kummer
normalization is smooth. Thus the special fibre consists
exactly of two rational components joined by $\ell$ nodes,
each of thickness $a$. It is stable because $\ell\ge3$;
no regularity of the total space at a node of higher thickness
is required.

The skeleton $\Gamma$ has two vertices and $\ell$ parallel
edges $e_1,\ldots,e_\ell$, all of length $a$.
Its node annuli $A_j$ are disjoint copies of
$\{|\alpha|<|T|<1\}$ with coordinate $T$.
This last all-sheet assertion can also be checked directly:
on that annulus
\[
 \frac{T(T-1)}{T-\alpha}
                =-\frac{1-T}{1-\alpha/T},
\]
and both principal units have convergent $\ell$-th roots.
The full inverse image is therefore $\ell$ copies of the annulus.

\emph{An integral character of slope one.}
Choose the rational branch point over $T=0$ as base point for
an algebraic Abel--Jacobi map $i:Z\to J$.
The rational components imply that the Jacobian is totally
degenerate. Over $\Omega$ its uniformization is
\[
                q:T^{\mathrm{an}}\longrightarrow J_\Omega^{\mathrm{an}},
                  \qquad\ker q=\Lambda.
\]
The uniformization is a topological covering and local analytic
isomorphism with free discrete period lattice, as recalled in
\citet[Sections~4.1--4.2]{D_BakerRabinoff2015}.
That source's Section~4.5.2 identifies the character lattice
with $M=H_1(\Gamma,\Z)$; Proposition~6.1 and Corollary~6.6
identify the tropicalized Abel--Jacobi map with integration on
$\Gamma$ using the edge-length pairing.

An open Berkovich annulus has an interval as skeleton and
open-ball branches; it is locally path connected and simply
connected. Thus $i|_{A_j}$ lifts continuously to the covering
torus, and the local analytic inverses of $q$ make this lift
an analytic map $b_j:A_j\to T^{\mathrm{an}}$.
Orient every edge by increasing $v(T)$. For $j'\ne j$, the
cycle $c_j=e_j-e_{j'}$ is an integral element of $M$.
Let $\chi_j$ be its character and put $u_j=\chi_j\circ b_j$.
Integration of the harmonic form with edge coefficients $c_j$
along a segment of $e_j$ equals the length of that segment.
Consequently
\[
                       \frac{d\,v(u_j)}{d\,v(T)}=1.
\]
This is character evaluation in the stated integral lattice,
not evaluation after inversion of a period matrix. Different
torus lifts add a constant tropical period and do not change
the derivative.

\emph{The middle collar and arbitrary root orders.}
The annular-unit factorization used in D.16 gives
\[
       u_j=c_j'T(1+h_j),\qquad
       h_j=\sum_{n\ne0}h_{j,n}T^n,\qquad c_j'\in\Omega^*.
\]
For clarity, the term $c_j'T$ is the actual degree-one Laurent
coefficient. At a Gauss radius in $|\Omega^*|$, the reductions
of a unit and its inverse are inverse units in the Laurent
polynomial ring, so exactly one monomial dominates.
The exponent cannot change inside the annulus: convergence
on compact subannuli reduces a proposed change to finitely
many terms, and a tie occurs at a radius in $|\Omega^*|$
because $v(\Omega^*)=\Q$. Such a tie is impossible.
The known slope forces the exponent to be one everywhere.

Dominance at all radii $|\alpha|<r<1$, followed by limits
at the endpoints, gives for $n>0$
\[
                         |h_{j,n}|\le1,
                   \qquad |h_{j,-n}|\le|\alpha|^n.
\]
It follows at $s=v(T)$ that
\begin{equation}\label{eq:D17:middle-bound}
 |h_j|\le p^{-\min\{s,a-s\}}<p^{-D}
                      \quad(D<s<a-D).
\end{equation}
Every closed subannulus of this collar has a strict uniform
margin. The logarithm and exponential therefore converge
locally uniformly, and
\[
                b_j'=\exp\bigl(m^{-1}\log(1+h_j)\bigr)
\]
is analytic with $(b_j')^m=1+h_j$.
Indeed $D>1/(p-1)$ gives
$|\log(1+h_j)|\le|h_j|$, and the exponent has valuation
strictly greater than $D-v_p(m)>1/(p-1)$.
These inequalities explicitly include $p=2$ and arbitrary
$p$-power divisors of $m$.

Now form the single algebraic finite \emph{\'etale} cover
\[
                    Z'=Z\times_{i,J,[m]}J\longrightarrow Z.
\]
It is \emph{\'etale} for every $m$ because the ground field
has characteristic zero. On the full inverse image of each
$A_j$, the period-lattice argument, applied to $a_A=b_j$,
gives an $m$-th root of $u_j$. Dividing by the pullback of
$b_j'$ and by an $m$-th root of $c_j'$ gives an $m$-th root
of $T$ over the middle collar. The $A_j$ are disjoint, so the
functions combine on the entire inverse image of $A^\circ$.
Choose any connected component $Q$ of $Z'$ over $\bar k$.
Its map to the connected $Z$ is finite \emph{\'etale} of
positive degree and therefore surjective. It is smooth and
proper, and retains the roots on every inverse-image sheet.
The composite $Q\to\mathbb P^1$ has only horizontal indices
$1$ and $\ell$, because its upper map to $Z$ is \emph{\'etale}.
All covers used are algebraic over $\bar k$; their analytic
roots and uniformization lifts need only exist over $\Omega$.
\end{proof}

\begin{claim}{D.18}{Proved}\label{clm:D18}
For the fixed prolongation of the real-valued rank-one norm in
D.15, the specified subfield $K^{\mathrm{sep}}$ of its completion
$E$ is contained in $H$, the closure of $\widetilde K$ inside
$E$. In particular $H=E$.
\end{claim}
\begin{proof}
We reconstruct R24, Section~6, with the field copies retained
throughout \citep{D_OpenAI2026}.
Let $F/K_0$ be any finite Galois extension inside the fixed
$K^{\mathrm{sep}}$, and let $C\to X_{\bar k}$ be its smooth
proper curve cover. Its branch set $B$ is finite. If it is
empty, then $F\subset\widetilde K\subset H$ by definition.
Suppose $B\ne\varnothing$.

For every $b\in B$, choose a rational local parameter $t_b$
and a disk $B_b$ as in D.15: the disks are pairwise disjoint,
avoid $\xi$, and the full cover above each is a disjoint union
of maps $t_b=z_c^{e_c}$ on disks. Put
$m_b=\operatorname{lcm}_{c\mid b}e_c$.
Choose $\beta_b\in\bar k^*$ whose absolute value is smaller
than the radius of $B_b$ in the coordinate $t_b$.
Choose an integer
$D_b>v_p(m_b)+1/(p-1)$ and
$\alpha_b\in\bar k^*$ with $a_b=v(\alpha_b)>2D_b$.
D.17 gives a cover of the $T$-line with horizontal indices
dividing $\ell$ and a full inverse-image root of $T$ on
$D_b<v(T)<a_b-D_b$.

Pull this cover back by the finite nonconstant morphism
$X_{\bar k}\to\mathbb P^1_{\bar k}$ defined by
$T=t_b/\beta_b$, normalize, and take a connected component
$W_b\to X_{\bar k}$.
The horizontal-index statement of D.15 shows that its indices
still divide $\ell$. The analytic roots pull back to every
point of this normalized component over the collar
\[
 |\beta_b|p^{-(a_b-D_b)}<|t_b|<|\beta_b|p^{-D_b}
                         \quad\text{inside }B_b.
\]
Multiplication by an $m_b$-th root of $\beta_b$ in $\Omega$
gives a root of $t_b$ there.

Choose embeddings of the finitely many fields $\bar k(W_b)$
into the fixed $K^{\mathrm{sep}}$ over $K_0$, and form their
compositum. Let $W\to X_{\bar k}$ be its smooth proper
connected curve. It dominates every $W_b$, so all the roots
persist on the full inverse collars in $W_\Omega^{\mathrm{an}}$.
D.15's compositum calculation keeps every horizontal index a
divisor of $\ell$; it uses the residue-characteristic-zero
closed-point valuations, not the vertical reduction valuations.

Apply the all-sheet patching construction of D.15 with
$Y=W_\Omega$ and the just-constructed collars. It gives a
connected finite \emph{\'etale} cover
$V_\Omega\to W_\Omega$ and an analytic map
\[
                  (V_\Omega^{\mathrm{an}})|_U
                          \longrightarrow C_\Omega^{\mathrm{an}}
                         \quad\text{over }X_\Omega^{\mathrm{an}},
\]
where $U$ is the exterior of the smaller closed branch disks
and $\xi\in U$.
Since $W$ is proper and $\bar k\subset\Omega$ are algebraically
closed fields, invariance of finite \emph{\'etale} covers under
this extension gives a connected finite \emph{\'etale} cover
$V\to W$ and an isomorphism
\[
                         V\times_{\bar k}\Omega\simeq V_\Omega
                                  \quad\text{over }W_\Omega
\]
\citep[Expos\'e~X, Corollary~1.8]{E_SGA1}.
Only this finite \emph{\'etale} cover is descended, not the
analytic exterior map or its root trivializations.

The composite $V\to X_{\bar k}$ has horizontal indices
still dividing $\ell$, because its upper map is \emph{\'etale}.
Extend the selected embedding of $\bar k(W)$ to
$\bar k(V)$ inside $K^{\mathrm{sep}}$.
Its Galois closure has the same index bound by D.15, so the
image of $\bar k(V)$ lies in $L_\ell$ and hence in $H$ by D.16.
This membership does not require $V\to X$ itself to be Galois.

The evaluation statement of D.15 now gives a $K_0$-embedding
$F=\bar k(C)\hookrightarrow H$. Explicitly, the chosen
embedding of $\bar k(V)$, together with $\Omega\subset H$,
gives an $H$-valued point of $V_\Omega$ over the originally
specified $H$-valued point of $X_\Omega$. Its analytic image
lies over $\xi\in U$. Evaluate the exterior analytic map
using the induced isometric map of completed residue fields.
The resulting point of $C_\Omega(H)$ projects to the generic
point of $X_{\bar k}$ with the specified inclusion $K_0\to H$.
Since $C\to X_{\bar k}$ is finite, its support on $C$ is
generic as well. It therefore gives the stated field embedding.
Keeping the entire $H$-valued point makes the argument valid
even if $\xi$ itself is an $\Omega$-rational analytic point.

Finally this embedding lands in the specified copy of $F$.
For each $x\in F$, its minimal polynomial over $K_0$ splits
in $F$ because $F/K_0$ is normal. Its image in $H\subset E$
is one of these same roots: a split polynomial over a field
cannot acquire additional roots in a larger field. Thus the
image is contained in $F$, and equality follows from
$[F:K_0]$. Consequently the chosen $F\subset E$ is contained
in $H$.

Every element of $K^{\mathrm{sep}}$ lies in such a finite
Galois extension of $K_0$, since $K^{\mathrm{sep}}$ is an
algebraic closure of $K_0$ in characteristic zero. Taking their
union proves $K^{\mathrm{sep}}\subset H$ inside $E$.
The subfield $H$ is closed in $E$ and contains its dense subfield
$K^{\mathrm{sep}}$, so $H=E$.
\end{proof}

\subsection{From the audited completion theorem to local sections}

\begin{claim}{D.19}{Proved; the section-theoretic consequence of R24}
Let $k/\Q_p$ be finite, and let $X/k$ be smooth, proper,
geometrically connected of genus at least two. For the full arithmetic
\'{e}tale group, the map
\[
 X(k)\longrightarrow
 \{s:G_k\longrightarrow\pi_1(X):s\text{ continuous},\;
                         \mathrm{pr}\circ s=1\}/\Pi_X
\]
is bijective. Here the quotient is by conjugation by the full geometric
group. We derive this from D.18, with no birational-lift hypothesis on $s$.
The argument is the section-theoretic argument of the supplied R24,
whose inputs and valuation comparisons we check explicitly below.
\end{claim}
\begin{proof}
Put $K=k(X)$ and let $\widetilde K\subset K^{\mathrm{sep}}$
be the extension corresponding to $\pi_1(X)$. We first show that,
for every rank-one valuation $w$ of $K$ extending the $p$-adic
valuation, and every chosen prolongation $w^{\mathrm{sep}}$,
restriction induces an isomorphism of profinite groups
\[
 D_{w^{\mathrm{sep}}}(K^{\mathrm{sep}}/K)
       \xrightarrow{\ \sim\ }
 D_{\widetilde w}(\widetilde K/K),
 \qquad \widetilde w=w^{\mathrm{sep}}|_{\widetilde K}.       \tag{D.19.1}
\]
A rank-one ordered value group embeds in $\R$; normalize $w(p)=1$.
For surjectivity, lift an element of the right decomposition group
to $G_K$. Its translate of $w^{\mathrm{sep}}$ is another prolongation
of $\widetilde w$. The kernel $G_{\widetilde K}$ acts transitively
on these prolongations, so an adjustment gives a lift fixing
$w^{\mathrm{sep}}$. The infinite transitivity assertion follows from
the finite normal-extension assertion and compactness of the Galois
group. An automorphism fixing this valuation preserves the normalized
absolute value: two order embeddings of a rank-one value group into
$\R$ differ by a positive scale, fixed by the value of $p$.
It therefore extends continuously to the completion
$E=\widehat{K^{\mathrm{sep}}}$. If it fixes $\widetilde K$, it fixes
its closure $H=\widehat{\widetilde K}\subset E$. D.18 gives the
specified containment $K^{\mathrm{sep}}\subset H$, so that
automorphism is the identity. Decomposition groups are closed:
preservation of membership in a valuation ring is a closed condition
on each finite Galois orbit of algebraic elements. Thus (D.19.1) is
a continuous bijection between compact Hausdorff groups and has a
continuous inverse.

We next prove that
\[
          D_{\widetilde w}(\widetilde K/K)\longrightarrow G_k
                \quad\hbox{has no continuous section}       \tag{D.19.2}
\]
for such a rank-one $w$. Otherwise (D.19.1) would produce a continuous
birational section $\widehat s:G_k\to G_K$ whose image fixes
$w^{\mathrm{sep}}$. The opened Proposition~2.4(2), together with
Definition~1.4 and Observation~1.3, of \citet{D_Koenigsmann2005}
says that, over a finite extension of $\Q_p$, every such birational
section for a smooth complete curve is contained in a decomposition
group of a rational closed-point valuation $\operatorname{ord}_a$,
for some $a\in X(k)$. Choose its prolongation fixed by this same image.

Let $\Sigma=\widehat s(G_k)$ and $M=(K^{\mathrm{sep}})^\Sigma$.
The subgroup is compact and closed, hence $G_M=\Sigma$. The
restrictions $v_1,v_2$ of the two chosen prolongations to $M$ are
henselian: their absolute Galois group fixes one prolongation and
acts transitively on all prolongations, so that prolongation is unique.
They remain nontrivial and rank one. Indeed value groups in an
algebraic valued extension have torsion quotient, which leaves the
chain of convex subgroups unchanged. They are inequivalent because
$v_1(p)=1$ whereas $v_2|_{k^\times}=0$.

We give the contradiction directly, including $p=2$ and nondiscrete
value groups. Choose a nonsquare $d\in k^\times$. It remains a
nonsquare in $M$: a square root would lie in the chosen $\bar k$
and be fixed by $\Sigma$, contradicting the surjection
$\Sigma\to G_k$. Choose a uniformizer $t$ of the local ring at
$a$, and an integer $N>v_1(t)$. Then $z=p^N/t$ satisfies
$v_1(z)>0$ and $v_2(z)<0$. For positive integers $r$ set
\[
 q_r=d+(1-d)\frac{z^r}{1+z^r}.
\]
The denominator is nonzero, since $v_1(z^r)>0$. Direct calculation gives
\begin{align*}
 v_1(q_r/d-1)&=v_1((1-d)/d)+r v_1(z)\longrightarrow+\infty,\\
 v_2(q_r-1)&=v_2(d-1)-r v_2(z)\longrightarrow+\infty.
\end{align*}
For large $r$ these exceed $2v_1(2)$ and $2v_2(2)$, respectively.
Hensel's lemma gives a square root of $q_r/d$ in $(M,v_1)$ and
a square root of $q_r$ in $(M,v_2)$. Their quotient makes $d$ a
square in $M$. This proves (D.19.2); it also supplies the exact
two-valuation argument instead of assuming a general uniqueness slogan.

Now take any continuous section $s:G_k\to\pi_1(X)$.
The opened Pop--Stix valuative localization theorem
\citep[Theorem~26 in the author preprint]{D_PopStix2017}
applies to a smooth hyperbolic geometrically connected curve over
a finite extension of $\Q_p$, with no assumed point or birational
lift. It supplies a valuation $w$ of $K$ extending the $p$-adic
valuation and a prolongation $\widetilde w$ to $\widetilde K$ such that
\[
                    s(G_k)\subset D_{\widetilde w}.
\]
The rational rank of $\Gamma_w/\Gamma_k$ is at most one.
For clarity, two values rationally independent modulo $\Gamma_k$
would be the values of elements $u,v\in K^\times$; distinct monomials
in any polynomial in $u,v$ over $k$ would have distinct valuations,
so no such polynomial could vanish. This would give transcendence
degree at least two, contrary to $K=k(X)$. Since $\Gamma_k$ has
rational rank one, the ordered rank of $w$ is at most two.

Rank one is excluded by (D.19.2). In rank two let $u$ be the
rank-one coarsening. The nonzero proper convex subgroup of
$\Gamma_w$ extends uniquely through an algebraic value-group
extension, by taking all elements a positive multiple of which
lies in that subgroup. Consequently $s(G_k)$ preserves the induced
prolongation of $u$. If $u$ is nontrivial on $k$, its restriction
is equivalent to the $p$-adic valuation, since that is the restriction
of $w$ and a nontrivial coarsening of a rank-one valuation is itself.
Normalization and (D.19.2) again exclude it.

Thus $u$ is trivial on $k$. Properness gives a center on $X$.
It is not the generic point because $u$ is nontrivial, hence is a
closed point $a$. Its valuation ring dominates the discrete valuation
ring $\mathcal O_{X,a}$. Writing each element of $K^\times$ as a
power of a uniformizer times a unit of that ring shows that $u$
is equivalent to $\operatorname{ord}_a$. The refinement $w$ is
the composition with the unique $p$-adic valuation on the finite
extension $k(a)/k$. A chosen horizontal prolongation selects a
geometric point over $a$, so its decomposition group maps to the
corresponding conjugate of $G_{k(a)}\subset G_k$. Since $s$ maps
onto $G_k$, this inclusion forces $k(a)=k$.

In the full \'{e}tale tower, the horizontal inertia at $a$ is
trivial. Its decomposition group maps isomorphically to $G_k$
and is precisely the point section, as is also recorded in
\citet[Section~6.1]{D_PopStix2017}. Therefore $s$ has the class
of $s_a$. Changes of geometric lift conjugate by $\Pi_X$,
which is exactly the ambiguity in the claim. This proves surjectivity.

For injectivity suppose $s_a$ and $s_b$ are geometrically conjugate.
Use $a$ to define $i_a:X\to J=\operatorname{Jac}(X)$ by
$x\mapsto[x-a]$. For every integer $n>0$, the map $[n]:J\to J$
is finite \'{e}tale, including $p\mid n$, because the field has
characteristic zero. The fibre at $i_a(a)=0$ has a fixed lift.
Conjugacy of sections preserves the existence of such a lift in
every finite \'{e}tale cover, hence
\[
                     i_a(b)\in\bigcap_{n\ge1}nJ(k).
\]
The compact $p$-adic analytic group $J(k)$ has a torsion-free open
subgroup $U\simeq\Z_p^{[k:\Q_p]\dim J}$ by the local logarithm.
For $q=[J(k):U]$ we have $qJ(k)\subset U$, so the displayed
intersection lies in $\bigcap_r p^rU=0$. Thus $[b-a]=0$.
If $a\ne b$, a rational function with divisor $b-a$ would have
degree one as a map to $\mathbb P^1$, contradicting positive genus.
Thus $a=b$ and the proof is complete.
\end{proof}

The argument verifies the local proper assertion of R24 through its
specified completion theorem. It supplies no almost-everywhere
integrality or global finite-descent equality over a number field;
the continuous global obligations D1--D4 remain separate.

\subsection{Two explicit obstructions to proposed global shortcuts}

The following refutations concern two proposed auxiliary inferences.
They do not refute D1--D4. In particular, a rational place and a rational
tangential parameter are different requirements.

\begin{claim}{D.20}{Refuted}\label{clm:D20}
The following implication is false: a continuous birational section over
$\mathbf Q$ whose tangential Kummer class is represented by an ordinary
element of $\mathbf Q_p^\times$ at every finite prime, and by a unit at a
set of primes of Dirichlet density one, must have a rational tangential
class or be a unit at all but finitely many primes. The failure occurs
for a section supported at a rational place.
\end{claim}
\begin{proof}
Choose increasing distinct primes $q_j>2^{j^2}$ and define
\[
 A=\prod_{j\geq1}q_j^{j!}\in
 \widehat{\mathbf Q^\times}
 =\varprojlim_n\mathbf Q^\times/(\mathbf Q^\times)^n.
\]
Modulo $n$th powers, every factor with $j\geq n$ is trivial. Thus the
product defines compatible finite-level classes. Its valuation
coordinates are
\[
 v_p(A)=\begin{cases}j!&p=q_j,\\0&p\notin\{q_j:j\geq1\}.
 \end{cases}
\]
Infinitely many are nonzero, so $A$ is not the image of an element of
$\mathbf Q^\times$.

The power completion of $\mathbf Q_p^\times$ is
\[
 \widehat{\mathbf Q_p^\times}
 \simeq p^{\widehat{\mathbf Z}}\times\mathbf Z_p^\times.
\]
Indeed, $\mathbf Q_p^\times=p^{\mathbf Z}\times\mathbf Z_p^\times$,
and the unit group is complete for its finite power quotients. Since
$v_p(A)$ is an ordinary integer, its local image has a unique actual
representative $a_p\in\mathbf Q_p^\times$. One may also construct $a_p$
directly. Remove the one possible factor with $q_j=p$; the remaining
product converges in $\mathbf Z_p^\times$, because $j!$ is eventually
divisible by the exponent of $(\mathbf Z/p^r\mathbf Z)^\times$ for each
$r$. Restore the removed factor. Hence $v_p(a_p)=v_p(A)$.

The set $\{q_j\}$ is infinite but has Dirichlet density zero:
\[
 \sum_j q_j^{-s}\leq\sum_j2^{-j^2}<\infty\qquad(s\geq1),
\]
whereas $\log(1/(s-1))$ tends to infinity as $s\downarrow1$.
The $a_p$ are therefore units at density one, but not cofinitely many,
primes. The real local class is the positive sign class.

We realize these data by a continuous section. Fix a prolongation of
the place $t=0$ of $\mathbf Q(t)$ and its tangential splitting
\[
 D_0\simeq\widehat{\mathbf Z}(1)\rtimes G_{\mathbf Q}
 \subset G_{\mathbf Q(t)}.
\]
For each $j$, choose compatible radicals of $q_j$ and the corresponding
continuous Kummer cocycle
$c_j:G_{\mathbf Q}\to\widehat{\mathbf Z}(1)$. The series
\[
 c(g)=\sum_{j\geq1}j!c_j(g)
\]
converges uniformly in $g$ modulo every $n$, because all terms with
$j\geq n$ vanish modulo $n$. Its finite-level cocycle identities imply
that it is a continuous cocycle. If $s_0$ denotes the standard
tangential section and $i$ the inertia inclusion, then
\[
 s_A(g)=i(c(g))s_0(g)
\]
is a continuous section of $G_{\mathbf Q(t)}\to G_{\mathbf Q}$.
Its pullback of the Kummer class of $t$ is $A$.

For completeness, write a rational function near zero as
$f=t^r u(t)$, where $u$ is regular and nonzero at zero. Hensel lifting
gives compatible roots of $u/u(0)$ in the strict henselian local field,
so the Kummer character of $u$ on $D_0$ is inflated from that of $u(0)$.
Consequently
\[
 s_A^*\kappa(f)=r[A]+\kappa(u(0)).
\]
Thus $s_A^*\kappa(t-c)=\kappa(-c)$ for every nonzero rational $c$,
whereas $s_A^*\kappa(t)=[A]$ is not rational. The underlying rational
place is still exactly zero. The example is cuspidal on
$\mathbf P^1-\{0,1,\infty\}$, and its extra tangential data disappear
in the proper unramified quotient, which kills this inertia. It is
therefore compatible with all the asserted rational-place conclusions.
\end{proof}

\begin{claim}{D.21}{Refuted}\label{clm:D21}
It is false that every fourth-power Fermat twist arising from a section
of $\mathbf P^1_{\mathbf Q}-\{0,1,\infty\}$ admits a nonconstant map
over $\mathbf Q$ to an elliptic curve of rank zero. The twist
\[
 C:\quad x^4+2y^4=3z^4
\]
arises above the rational point $t=1/3$, and every elliptic target of
$C$ over $\mathbf Q$ has positive Mordell--Weil rank.
\end{claim}
\begin{proof}
The plane quartic is smooth, since its three partial derivatives cannot
vanish simultaneously at a projective point; its genus is three. Its
open subset $xyz\ne0$ maps to the thrice-punctured line by
\[
 t=\frac{x^4}{3z^4},\qquad 1-t=\frac{2y^4}{3z^4}.
\]
The rational point $P=(1:1:1)$ lies above $t=1/3$, so this is already
a twist belonging to a geometric section.

The three sign involutions have genus-one quotients. For example, the
$x$-sign quotient has affine equation $v^2=3-2u^4$, where
$u=y/z$ and $v=(x/z)^2$. Its four distinct branch points give genus
one, and the image of $P$ is rational. The other two quotients have the
same description after changing coordinates. Explicit elliptic targets
and maps are
\[
\begin{array}{c|c|c|c}
 &\text{elliptic equation}&X&Y\\\hline
 E_x&Y^2=X^3-6X&-2y^2/z^2&-2x^2y/z^3\\
 E_y&Y^2=X^3-12X&-2x^2/z^2&-4xy^2/z^3\\
 E_z&Y^2=X^3+18X&3x^2/y^2&9xz^2/y^3.
\end{array}
\]
These rational maps extend to the smooth projective curves. To verify
their equations, put $F=x^4+2y^4-3z^4$. After clearing denominators,
$Y^2-X^3-AX$ is respectively $4y^2F$, $8x^2F$, and $-27x^2F$.
Each map factors through its indicated sign quotient. In the quartic
model $v^2=A u^4+B$, the factor map is
$X=A u^2$, $Y=Auv$, with equation $Y^2=X^3+ABX$; its degree is two.
After translating the image of a rational origin, this is an isogeny
of genus-one Jacobians. Thus $E_x,E_y,E_z$ represent the three quotient
isogeny classes.

The images of $P$ are
\[
 (-2,-2)\in E_x(\mathbf Q),\quad
 (-2,-4)\in E_y(\mathbf Q),\quad
 (3,9)\in E_z(\mathbf Q).
\]
They have infinite order. Exhaustive finite-field counts give
\[
\begin{array}{c|cc}
 &\text{first good-reduction count}&\text{second good-reduction count}\\\hline
 E_x&\#E_x(\mathbf F_5)=8&\#E_x(\mathbf F_{13})=18\\
 E_y&\#E_y(\mathbf F_5)=10&\#E_y(\mathbf F_7)=8\\
 E_z&\#E_z(\mathbf F_5)=10&\#E_z(\mathbf F_7)=8.
\end{array}
\]
The discriminant of $Y^2=X^3+AX$ is $-64A^3$, so all displayed
primes are good. Prime-to-$p$ rational torsion injects under reduction
at $p$; see \cite[Chapter~II, Proposition~5.4 and Corollary~5.7]{D_Milne2021}.
For $E_x$ the count at $13$ excludes $5$-primary torsion, and the
count at $5$ excludes $13$-primary torsion. All other primary torsion
orders divide $\gcd(8,18)=2$. The analogous argument at $5,7$ for
$E_y,E_z$ uses $\gcd(10,8)=2$. Each displayed point has nonzero
$Y$-coordinate, so is neither the identity nor of order two. Each
quotient isogeny class has positive rank. The exhaustive calculation is
reproducible from \path{code/D_Fermat_quartic_certificate.py}; its
output is \path{audit/D_Fermat_quartic_certificate.json}.

Finally, let $V$ be the Klein four group of projective sign changes,
with nonidentity elements $\tau_x,\tau_y,\tau_z$. The full quotient
has function field $\mathbf Q(x^2/z^2,y^2/z^2)$ with equation
$X^2+2Y^2=3$. This conic has the rational point $(1,1)$, and hence has
genus zero. On $J=\operatorname{Jac}(C)$ the full norm
$1+\tau_x+\tau_y+\tau_z$ is therefore zero, by pushforward and pullback
for $C\to C/V$. In $\operatorname{End}(J)\otimes\mathbf Q$ the
idempotents
\[
 e_i=\frac{1+\tau_i}{2}\qquad(i=x,y,z)
\]
are pairwise orthogonal and sum to one. Their images are the three
genus-one quotient Jacobians up to isogeny. Hence $J$ is
$\mathbf Q$-isogenous to their product. If $C$ maps nonconstantly to
an elliptic curve $E/\mathbf Q$, translate the image of $P$ to the
origin. The induced map $J\to E$ is nonzero, and so is its restriction
to at least one factor of this product up to isogeny. A nonzero
homomorphism between elliptic curves is an isogeny. Therefore $E$ is
isogenous to one of $E_x,E_y,E_z$ and has positive rank.

The same decomposition excludes a nonconstant map from $C$ to any
abelian variety $A/\mathbf Q$ with $A(\mathbf Q)$ finite: the induced
nonzero homomorphism $J\to A$ is nonzero on some elliptic factor, and
that factor's infinite rational points have infinite image because
its kernel is finite. In particular $C$ cannot map nonconstantly over
$\mathbf Q$ to a positive-genus standard $X_0(N)$ or $X_1(N)$.
Indeed, the proof of Stoll's Corollary~8.8 \cite{D_Stoll2007} supplies
a nonconstant map from each such modular curve to a modular abelian
variety of analytic rank zero, whose rational points are finite.
A nonconstant map of proper curves $C\to M$ is surjective, so its
composition with that map from $M$ would remain nonconstant, a
contradiction. Thus this particular twist does not meet the target
hypotheses of D.22 below.

This excludes a universal rank-zero elliptic-target shortcut on these
twists. It does not exclude additional covers or a different descent
argument, and makes no assertion about the exact ranks or
Tate--Shafarevich groups.
\end{proof}

\paragraph{Why finite localization is not yet finite support.}
Betts--Stix prove that, for a number field containing no CM subfield
and a smooth projective curve of genus at least two, the image of its
Selmer section set in the points over each fixed finite completion is
finite \cite[Theorem~A]{D_BettsStix2025}. The local theorem D.19 supplies
the required point-theoretic localizations at finite places. At a real
place one must separately use the classical real section theorem,
whose target is the set of connected components of the real locus;
the precise input is \cite[Corollary~3.13]{D_Mochizuki2003Real}.
D.19 itself is a $p$-adic theorem and does not address real places.
At a complex place the local Galois group is trivial. Subject to these
archimedean facts, proper global sections are Selmer and the stated
finite-localization theorem applies.

Its conclusion is not the common finite arithmetic support used by
Stix's corollary \cite[Section~3, Corollary~5]{D_Stix2015}. The finite
sets at different completions are not asserted to lie in the base
changes of a single finite subscheme over the number field. Even a
singleton local point can have generic scheme-theoretic image over
the number field if its coordinates are transcendental over that
field. The missing common-support assertion, or the cofinite
integrality condition of D.5, must therefore be supplied by an
additional global argument.

\begin{claim}{D.22}{Proved}\label{clm:D22}
Let $V/k$ be a smooth geometrically connected hyperbolic curve over a
number field, with smooth projective completion $C$. If $C$ admits a
nonconstant morphism to an abelian variety $A/k$ for which $A(k)$ is
finite and $\Sha(k,A)_{\mathrm{div}}=0$, then every continuous section
of $\pi_1(V)\to G_k$ is a rational-point section or is cuspidal.

In particular this conclusion holds over $\mathbf Q$ whenever $C$
admits a nonconstant $\mathbf Q$-morphism to a standard modular curve
$X_0(N)$ or $X_1(N)$ of positive genus. It includes the proper modular
curves of genus at least two and hyperbolic open modular curves with
positive-genus completions.
\end{claim}
\begin{proof}
Let $s$ be a continuous section. At every finite place $v$, D.6, using
the proper local theorem D.19, makes $s_v$ a point or cusp section.
Its image in the fundamental group of $C$ is therefore induced by a
point $x_v\in C(k_v)$. At every real place, the profinite real section
theorem \cite[Corollary~3.13]{D_Mochizuki2003Real} identifies sections
for the hyperbolic curve $V$ with connected components of $V(\mathbf R)$.
Choose a point in that component and map it to $C$. This specifies the
required real component of the modified adelic point; it does not
assert that a real point within a component is uniquely determined.
Complex factors are omitted in the corrected modified adelic space
\cite[Section~1]{D_StollErrata2017}. Thus $s$ supplies
\[
 (x_v)\in C(\mathbf A_k)_\bullet^{\mathrm{f\text{-}cov}}
 \subseteq C(\mathbf A_k)_\bullet^{\mathrm{f\text{-}ab}}.
\]
Here the full finite-descent assertion follows from Harari--Stix,
Theorem~2.1 and Remark~2.2(1) in the supplied preprint
\cite{D_HarariStix2012}: each finite torsor's evaluations are obtained
by restriction of its pullback along the global section. This is a
condition on torsors under arbitrary finite group schemes, not only
constant or geometrically connected torsors.

Under the stated abelian-target hypothesis, Stoll's Theorem~8.6
\cite{D_Stoll2007} gives
\[
 C(\mathbf A_k)_\bullet^{\mathrm{f\text{-}ab}}=C(k).
\]
Consequently the finite-place local points of $s$ all equal one
$x\in C(k)$, and their real components equal those of $x$. Choosing
$x$ as the real representative gives common finite support $\{x\}$.
If $s$ is cuspidal there is nothing further to prove. Otherwise Stix's
finite-support criterion, Corollary~5 in \cite{D_Stix2015}, shows
that $s$ is the section of a point of $V(k)$; its image in $C$ is
necessarily $x$. For a proper hyperbolic curve this is the point-section
conclusion, since no cusps occur.

For the modular assertion put $M=X_0(N)$ or $X_1(N)$, with
$g(M)>0$, and let $f:C\to M$ be nonconstant. Stoll's Corollary~8.8
states precisely that
$M(\mathbf A_{\mathbf Q})_\bullet^{\mathrm{f\text{-}ab}}=M(\mathbf Q)$.
Functoriality sends $(x_v)$ to one $y\in M(\mathbf Q)$.
The scheme $Z=f^{-1}(y)$ is finite over $\mathbf Q$, and
\[
 x_v\in Z(\mathbf Q_v)\qquad\text{for every finite place }v.
\]
Since $g(C)\geq g(M)>0$, Stoll's Theorem~8.2 applies to $Z\subset C$
and makes $(x_v)$ a diagonal rational point of $C$. This is the
finite-fibre proof of Proposition~8.5, with exactly the repair in
\cite[Section~3]{D_StollErrata2017}: no assertion that the specified
real component meets $Z(\mathbf R)$ is needed before applying
Theorem~8.2; a density-one set of finite places suffices. The
finite-support argument of the preceding paragraph now applies.
\end{proof}

\paragraph{Consequence for selected lifts of sections of $U$.}
If a section of $U=\mathbf P^1_{\mathbf Q}-\{0,1,\infty\}$ lifts to a
finite \'etale cover $V\to U$ whose completion satisfies D.22, its lift
is a point or cusp section. A finite \'etale map pushes such sections
to point or cusp sections, respectively: in the finite map between
the smooth completions the complement of $V$ lies over the complement
of $U$. This settles the original section in that case. Arbitrary
arithmetic twists of covers of $U$ are not thereby shown to map over
$\mathbf Q$ to a standard positive-genus modular curve or to a
rank-zero abelian target. Establishing that property for the remaining
twist classes would be an additional theorem. D.22 does not assert
universal D3 or D4.


\subsection{The exact unresolved obligations}

\begin{claim}{D.U1}{Unresolved: D1}
For every continuous section $s:G_\Q\to\pi_1(U)$ with
$C_{\Pi_U}(s(G_\Q^\circ))=1$, there exists $x\in U(\Q)$
such that $s$ is geometrically conjugate to $s_x$.
\end{claim}
We checked the local-to-global Kummer route, finite-support descent,
and the stronger birational criterion. None establishes uniform
adelicity or birational liftability for an arbitrary section.
The proved local result D.19 and its affine consequence D.6 give
local points or cusps without this missing global condition.
No conclusion about all Fermat twists follows
from a computation for one twist with a rank-zero quotient.

\begin{claim}{D.U2}{Unresolved: D2}
Every continuous splitting $G_\Q\to G_{\Q(t)}$ belongs to the
decomposition group of a rational place.
\end{claim}
D.3 settles the proposed arithmetic criterion for a given Kummer
class; D.4 corrects its evaluation formula. What remains is to
prove the common finiteness of the nonunit places needed for that
criterion, or another condition forcing the place. Density-one
integrality alone does not do so.

\begin{claim}{D.U3}{Unresolved: D3}
Every section of every hyperbolic curve over every field finitely
generated over $\Q$ admits a birational lift.
\end{claim}
D.1 gives the exact equivalence with the section conjecture. No
solution of the required global finite embedding problems, with
compatible choices, has been constructed here. An abelianized
$H^2$ obstruction supplies a necessary condition only; its
vanishing is not a splitting of the nonabelian extension.

\begin{claim}{D.U4}{Unresolved: D4}
Every continuous section of $\pi_1(U)\to G_\Q$, for
$U=\mathbb P^1_\Q\setminus\{0,1,\infty\}$, lifts through
$G_{\Q(t)}\to\pi_1(U)$.
\end{claim}
D.2 proves exactly what combining this with D2 would yield.
The remaining kernel is generated by the inertia of all deleted
closed points; replacing it by an abelian product does not solve
the original extension problem.

